\documentclass[11pt,a4paper]{article}
\usepackage[T1]{fontenc}
\usepackage[utf8]{inputenc}
\usepackage{lmodern}
\usepackage[margin=27mm,headheight=15pt]{geometry}
\usepackage{amsmath,amssymb,amsthm,mathtools}
\usepackage{microtype,booktabs,array,longtable}
\usepackage[dvipsnames]{xcolor}
\usepackage{enumitem,listings,fancyhdr}
\usepackage[unicode,pdfencoding=auto,psdextra]{hyperref}
\usepackage{xurl}% write: line breaks in long paths and labels
\hypersetup{colorlinks=true,linkcolor=MidnightBlue,citecolor=MidnightBlue,urlcolor=MidnightBlue,
 pdftitle={Smooth Everywhere, Universal on the Integers},
 pdfauthor={Research prepared for Vladimir Reshetnikov with ChatGPT},
 pdfsubject={Canonical quartic Diophantine certificates, arithmetic smoothness, and local solubility}}
\setlist{itemsep=3pt,topsep=5pt}
\setlength{\parskip}{4pt}
\setlength{\parindent}{1.1em}
\emergencystretch=2em
\pagestyle{fancy}
\fancyhf{}
\fancyhead[L]{\small Smooth Everywhere, Universal on the Integers}
\fancyhead[R]{\small ProveIt research}
\fancyfoot[C]{\thepage}
\renewcommand{\headrulewidth}{0.3pt}
\newtheorem{theorem}{Theorem}[section]
\newtheorem{lemma}[theorem]{Lemma}
\newtheorem{proposition}[theorem]{Proposition}
\newtheorem{corollary}[theorem]{Corollary}
\theoremstyle{definition}
\newtheorem{definition}[theorem]{Definition}
\newtheorem{example}[theorem]{Example}
\newtheorem{question}[theorem]{Research question}
\theoremstyle{remark}
\newtheorem{remark}[theorem]{Remark}
\newcommand{\N}{\mathbb N}
\newcommand{\Z}{\mathbb Z}
\newcommand{\Q}{\mathbb Q}
\newcommand{\R}{\mathbb R}
\newcommand{\A}{\mathbb A}
\newcommand{\Fp}{\mathbb F_p}
\newcommand{\Zp}{\mathbb Z_p}
\newcommand{\Spec}{\operatorname{Spec}}
\newcommand{\zeros}[2]{Z_{#1}(#2)}
\newcommand{\dd}{\mathrm d}
\newcommand{\height}{\operatorname{ht}}
\newcommand{\supp}{\operatorname{supp}}
\newcommand{\ann}{\operatorname{Ann}}
\newcommand{\guard}{g}
\newcommand{\source}{\mathcal S}
\newcommand{\variety}{\mathcal X}
% Write phase (batch 79, cluster J1): marks text added when the manuscript
% was written into the collection; unmarked text is the source's own.
\newcommand{\wtag}{\textup{[write]}}
\newcommand{\Dplus}{\mathcal D^{+}}
\lstset{basicstyle=\ttfamily\small,breaklines=true,columns=fullflexible,
 frame=single,rulecolor=\color{black!25},showstringspaces=false,
 xleftmargin=5pt,xrightmargin=5pt,aboveskip=8pt,belowskip=8pt}

\begin{document}
\hypersetup{pageanchor=false}
\begin{titlepage}
\thispagestyle{empty}
\vspace*{2mm}% write: was 16mm, to keep the added status sentence on this page
{\Large\color{MidnightBlue} PROVEIT\quad /\quad RESEARCH MANUSCRIPT\par}
\vspace{12mm}
{\Huge\bfseries Smooth Everywhere,\\[4pt]
Universal on the Integers\par}
\vspace{7mm}
{\Large Canonical quartic certificates with five extra variables,\\
no local obstruction, and explicit Jacobian identities\par}
\vspace{12mm}
{\large Prepared for Vladimir Reshetnikov\par}
{\normalsize Developed with ChatGPT\par}
{\normalsize 2 October 2026\par}
\vspace{10mm}
\noindent\rule{\textwidth}{0.5pt}
\begin{abstract}
We give a uniform arithmetic-geometric finalizer for Diophantine computation.
Given any finite system of quadratic integer equations in $n$ variables, it
constructs one polynomial of exact degree four in $n+5$ variables. Its natural
zeros are in explicit bijection with the original natural zeros, and its integer
zeros are exactly two sign-related copies of the original integer zeros.
Nevertheless, its affine scheme is smooth over $\Z$, every geometric fibre is
integral, a nonnegative dyadic rational point can always be constructed, and
there are integral $p$-adic points for every prime. Rational points are dense in
the real locus, which has exactly two path components. A degree-bounded
integral B\'ezout identity certifies smoothness without a genericity search.

The construction yields a sharply specified undecidable class of smooth,
everywhere locally soluble quartic equations. It also preserves canonical
bounded-execution certificates without increasing witness multiplicity. We
supply a counter-machine frontend, a complete worked export, a sharp
weighted-guard obstruction in this architecture, and
21,128 successful symbolic and finite checks. General undecidability for
smooth varieties is known; the contribution proposed here is the explicit
normal form and its combined guarantees, not a claim of priority for that
general phenomenon.
\end{abstract}
\vfill
\noindent\textbf{Status.} Ordinary mathematical proofs and executable exact checks.
Not a Lean/Coq formalization; not peer reviewed. No finite-fold MRDP theorem,
no solution to Hilbert's tenth problem over $\Q$, and no improvement to the
repository's universal operation-count records is claimed.

\smallskip\noindent{\small\raggedright\wtag\ Written into ProveIt's research-report collection
(\emph{hilbert-tenth-problem}) from batch-79 manuscript~13. ``Smooth'' means
the affine model is smooth over~$\Z$; Lemma~\ref{sdf:lem:lifting} and
Theorem~\ref{sdf:thm:counter} re-present results of the collection's report
\emph{canonical-diophantine-certificates}. See the
\hyperref[sdf:sec:preface]{editorial preface} and
Appendix~\ref{sdf:app:provenance}.\par}
\end{titlepage}
\hypersetup{pageanchor=true}

\tableofcontents
\newpage

\phantomsection\label{sdf:sec:preface}
\addcontentsline{toc}{section}{Editorial preface}
\section*{Editorial preface}

\wtag\ This preface, the bracketed notes marked \wtag\ in the text,
Appendix~\ref{sdf:app:provenance} and the bibliography entries marked \wtag\
were written when the manuscript was placed in ProveIt's research-report
collection. Everything else is the manuscript as delivered, with its
numbering of sections, statements and equations. Its labels gained the
prefix \texttt{sdf:}, and labels were added to three numbered results that
had none (Theorem~\ref{sdf:thm:weights}, Corollaries~\ref{sdf:cor:five}
and~\ref{sdf:cor:nobound}), to some sections, to the research questions and
to Appendix~\ref{sdf:app:export}.

\paragraph{The source.} The report is built from one manuscript, batch-79
manuscript~13 of ProveIt's incoming reports (archive
\path{Smooth_Quartic_Diophantine_Certificates.zip}, arrival commit
\texttt{ef2fc7990}, placed by commit \texttt{224ca41df}). It names the
repository commit \texttt{a845ab5d0} as the snapshot it compared against
(Section~\ref{sdf:sec:repo} and Appendix~\ref{sdf:app:provenance}). It is an
AI-assisted research manuscript (``Developed with ChatGPT'') prepared for
Vladimir Reshetnikov. Its proofs are conventional and its checks are finite
exact computations; it is unrefereed, and nothing in it is formalized in
Lean or Rocq.

\paragraph{Already reviewed.} After placement, the Hilbert's-tenth research
tree reviewed the archive (commit \texttt{fd4a2e8d3}; files
\path{review_smooth_quartic_aebfa386e.md}, \texttt{.py} and \texttt{.json}
in \path{Computability/HilbertTenthProblem/Papers/research-wip/native-stream-queue/},
indexed in \path{incoming_substrate_review_aebfa386e.md})~\cite{review13}.
Its verdict is ``PASS within the stated mathematical and implementation
scope'': it replayed all 21,128 assertions of the delivered verifier and
reproduced both example exports byte for byte, checked the integral Jacobian
identity in every characteristic, the fibre bijections, geometric
integrality, the dyadic and $2$-adic constructions and the restricted
guard-weight optimality, and ran independent checks (12,636 integer and 640
natural fibre tuples, 120 counter runs against a separate interpreter). It
found no theorem defect and no required repair. It stresses five scope
points that the text below already states: the size bounds of
Proposition~\ref{sdf:prop:size} count monomials, not arithmetic operations;
``degree four'' is witness degree, and a loader parameter treated as a
coordinate gives total degree six (Section~\ref{sdf:sec:parameter-degree});
the counter frontend has an external horizon and accepts halting \emph{by}
that horizon, so it is not a fixed-arity compiler; fixed-arity universality
comes only from an existing MRDP relation and proves no single-fold MRDP
theorem; and the code is not an authenticated checker for edited certificate
objects or hostile JSON. The finalizer gives no smaller operation schedule
for the repository's 87-operation universal polynomial.

\paragraph{What is new and what is re-proved.} The main new content is the
arithmetic geometry of the final equation: smoothness of the affine model
over~$\Z$ with an explicit integral B\'ezout certificate
(Sections~\ref{sdf:sec:smooth}--\ref{sdf:sec:integrality}), local and dyadic
points (Section~\ref{sdf:sec:local}), rational density and the two real
components (Section~\ref{sdf:sec:real}), the weighted-guard classification
(Section~\ref{sdf:sec:guards}) and the power-of-two padding
(Section~\ref{sdf:sec:power-two}). Three pieces re-present results already in
the collection's report \emph{canonical-diophantine-certificates} (CDC
below)~\cite{cdcreport}; they are printed because the argument uses them, with
\wtag\ notes at each and no novelty claimed:
\begin{itemize}
\item Lemma~\ref{sdf:lem:lifting}, natural version, is CDC's unique-extension
quadratization, \path{cdc:wf:lem:quadratization}, with the same proof
($F^+-F^-$ circuits, one natural variable per gate, uniqueness by induction
along the circuit); the integer version with subtraction gates is the same
induction. The sum-of-squares step $H=\sum q_i^2$ is CDC's
\path{cdc:wf:eq:sos}.
\item Theorem~\ref{sdf:thm:counter} is a second presentation of CDC's
counter-program adapter: one-hot control equations, the zero guard
\eqref{sdf:eq:zeroedge}, which is CDC's \path{cdc:eq:zero-guard}, the
externally fixed horizon of CDC's Section \path{cdc:pt:sec:boundary}, and
the halt self-loop caveat of \path{cdc:ct:sec:substrates}. New here are
only the exact counts $n_T$, $m_T$ and the positive-edge slack equation
\eqref{sdf:eq:posedge} with natural selectors summing to one in place of
Boolean rows.
\item The computability core of Theorem~\ref{sdf:thm:hardness} is CDC's
quartic normalization, \path{cdc:of:prop:quartic} (every computably
enumerable set has a quartic representation by MRDP, circuit lifting and a
sum of squares), together with Poonen's smooth-variety undecidability, which
the manuscript credits itself. The new content of that theorem is the
geometric package (a)--(e).
\end{itemize}
Three further credits are added as \wtag\ notes: Ramanujan's list of
universal diagonal quaternary forms for Lemma~\ref{sdf:lem:1122}, Minsky's
two-counter universality beside the formalized reductions of
Section~\ref{sdf:sec:counter}, and a variable count beside Poonen's lemma in
Section~\ref{sdf:sec:novelty}.

\paragraph{A non-claim of the collection met in finalized form.} The report
\emph{probabilistic-quantum-and-continuous-computation} (PQC)~\cite{pqcreport}
proves that each zero of its quartic Boolean-circuit energy $\mathcal E_C$
is a nondegenerate minimum (\path{pqc:ql:thm:hessian}) and, in the remark
``Nondegenerate minimum versus smooth hypersurface'' that follows it, does
not claim that the hypersurface $\mathcal E_C=0$ is nonsingular. It cannot
be: $\mathcal E_C$ is a sum of squares, so its gradient vanishes at every
real zero, which is the observation that opens Section~1 below. The residuals
of $\mathcal E_C$ are integer polynomials of degree at most two, so
Theorem~\ref{sdf:thm:main} applies to them. Five further variables give a
quartic whose affine hypersurface is smooth over~$\Z$, with the same natural
(here Boolean) zeros. The price is positivity. The finalized polynomial
takes negative real values, so its zeros are no longer energy minima, and the
Morse statement of PQC and the smoothness statement here concern different
polynomials. Neither report proves the other's statement.

\paragraph{Relation to CDC Part~XI.} CDC's Part~XI classifies
representations by polynomials that are nonnegative on the whole real orthant
(the classes $\Dplus_d$ of \path{cdc:of:thm:classification}). The finalized
polynomial $F$ of \eqref{sdf:eq:F} is never in such a class: at the dyadic
point \eqref{sdf:eq:dyadic} the coordinates $t,u,z_j$ are positive, and
$F_u=2tR=2B/K^3$ and $F_{z_j}=\lambda_j(4z_j-1)=2\lambda_jA_j/K^2$ are not
all zero, because $A_1^2+A_2^2+2A_3^2+2B^2=K^4-2c>0$. Decreasing one
coordinate with a nonzero derivative slightly gives a point with nonnegative
coordinates where $F<0$. Some sign change is forced: a polynomial that is
nonnegative near one of its zeros has zero gradient there, so a hypersurface
smooth at a zero inside the open orthant changes sign there. The manuscript
uses positivity only on the integer lattice (Lemma~\ref{sdf:lem:guard}),
which is among the settings that CDC's README, in its scope note for
manuscript~09 (CDC's Part~XI), says that Part's theorem does not cover
(``positivity only on natural points'').

\paragraph{Batch-79 siblings.} Several manuscripts of the same batch end with
an ordinary sum of squares of integer residuals of degree at most two in
natural unknowns, which is exactly the input of Theorem~\ref{sdf:thm:main}:
manuscript~03's bounded-bit quartic (placed for CDC's Part~XVII), and the
bounded ordinary exports of manuscripts~06 and~15 (``Coefficient extraction
at an external bound'' and ``Bounded ordinary quartic export'', Parts~I
and~II of the report \emph{polynomial-witness-histories}). Each could be
followed by this finalizer at the cost of five variables; no source states
that composition, and none cites this manuscript, which arrived after four of
them. Manuscript~02's exports keep fixed-base power atoms and are not
ordinary polynomial equations, so the finalizer does not apply to them as
they stand; the unbounded exports of~06 and~15 have polynomial-valued unknowns
and are outside its integer setting. Manuscript~15's question whether one
equation can have degree two or three is not answered here: the finalizer
stays at degree four.

\paragraph{Notation.} The manuscript reuses several letters between its
algebraic sections and its counter frontend, and some coincide with letters
of neighbouring reports. No symbol was renamed. The table lists the readings;
the right column gives the other reading in this report, or the
tempting false one.

\begin{center}
\small
\begin{longtable}{@{}>{\raggedright\arraybackslash}p{0.15\textwidth}%
 >{\raggedright\arraybackslash}p{0.38\textwidth}%
 >{\raggedright\arraybackslash}p{0.39\textwidth}@{}}
\toprule
Symbol or term & Main meaning & Other meaning in this report, or the tempting false reading\\
\midrule
\endhead
$F$ & the finalized quartic \eqref{sdf:eq:F} in $n+5$ variables & not nonnegative over $\R$ or on the orthant; not CDC's MRDP polynomial $F(x,y)$ nor PQC's energy $\mathcal E_C$\\
$g$ & the guard $\guard(z)=2z^2-z$ & Section~\ref{sdf:sec:counter}: $g$ is the number of guarded edges\\
$E$ & the Euler combination $E=4H$ (\eqref{sdf:eq:EDJ}, Section~\ref{sdf:sec:guards}) & Section~\ref{sdf:sec:counter}: the number of edges\\
$D_j$ & $D_j=F_{z_j}$ (Section~\ref{sdf:sec:smooth}) & countdown example: $D$ is the decrement edge\\
$H$ & $\sum_i q_i^2$ & countdown example: the halt self-loop edge\\
$t$, $q_i$, $c$ & homogenizing coordinate; homogeneous lift; $c=\sum_if_i(0)^2$ \eqref{sdf:eq:c} & Section~\ref{sdf:sec:counter}: target state $t(e)$, state variable $q_{s,i}$, counter $c_{s,j}$\\
$p$, $d$, $h$ & a prime; the power-of-two degree of Theorem~\ref{sdf:thm:padding}; $h(z_1)$ of \eqref{sdf:eq:2adic} & Section~\ref{sdf:sec:counter}: number of positive-decrement edges, slack $d_{s,e}$, halt state $h$\\
$B$ & the coordinate with $K^4-2c=A_1^2+A_2^2+2A_3^2+2B^2$ (Theorem~\ref{sdf:thm:dyadic}) & Corollary~\ref{sdf:cor:height}: the box radius\\
$w_j$, $W$ & guard weights and their total (Section~\ref{sdf:sec:guards}) & $w$: a gate variable (Lemma~\ref{sdf:lem:lifting}) and the second guard variable of \eqref{sdf:eq:bad-weighted}; $w_j$: rescaled guard coordinates in the proof of Theorem~\ref{sdf:thm:components}\\
$N$ & the integer $K^4-2c$ (Section~\ref{sdf:sec:local}) & not $\N=\{0,1,2,\ldots\}$ (as in CDC)\\
smooth, good reduction & the affine scheme $\variety\to\Spec\Z$ is smooth & not smoothness of a projective closure, nor a smooth proper model (neither claimed)\\
dyadic point & a point of $\variety(\Z[1/2])$, here with nonnegative coordinates & not density of dyadic points (not claimed)\\
degree four & degree in the witness variables of a specialized equation & not total degree with input parameters as coordinates (six, Section~\ref{sdf:sec:parameter-degree})\\
canonical, witness-faithful & natural zeros in bijection with the source's & not single-fold or finite-fold MRDP (not claimed)\\
finalizer & a post-processing of a quadratic system into one equation & not the operation-count ``finalizers'' of the Hilbert's-tenth research tree; this one costs more operations than a sum of squares\\
\bottomrule
\end{longtable}
\end{center}

\section{The problem: preserve computation, simplify the geometry}\label{sdf:sec:intro}

A computational certificate normally begins as a list of equations. States,
registers, scheduling choices, memory contents, and auxiliary multiplication
gates are constrained by polynomial residuals. Combining those residuals by
squaring is convenient over the integers, but it destroys geometric regularity:
if
\[
 P(x)=\sum_{i=1}^m f_i(x)^2,
\]
then every common zero of the $f_i$ is a critical point of $P$. This is not an
incidental implementation defect. The ordinary sum-of-squares finalizer makes
precisely the computationally meaningful points singular.

The question considered here is more demanding than whether computation has
some Diophantine representation:
\begin{quote}
Can a certificate be replaced by one low-degree equation, without changing
its natural-number witnesses, while making the affine model smooth at every
prime and making all local and rational existence tests automatically succeed?
\end{quote}
The answer is affirmative. The resulting separation is between lattice
arithmetic and geometry: smoothness, rational points, real connectedness, and
local integral points can all be harmless while integer solvability still
carries an arbitrary computably enumerable problem.

\subsection{Relation to the repository}\label{sdf:sec:repo}

The ProveIt repository already has a substantial collection of witness-faithful
Diophantine frontends. The reviewed canonical-certificate collection includes
bounded causal traces, queue and memory histories, dynamic heaps, interaction
nets, and stabilization certificates. Its status document explicitly separates
ordinary proofs, finite Python checks, and unproved uniform finite-fold claims
\cite{repo,canonical,status}. This manuscript adds a different layer: a
\emph{finalizer} that can follow any quadratic frontend, irrespective of its
underlying substrate.

The distinction matters. Replacing a scheduler, changing a memory encoding,
or quotienting commuting histories is a semantic task. Our finalizer neither
performs nor assumes such a quotient. It takes whatever witness set has already
been justified and preserves it exactly over $\N$. Thus a unique certificate
remains unique, a family indexed by trace classes retains those classes, and
an existing noncanonical representation remains noncanonical.

The targeted repository comparison used the canonical-certificate status file
at commit
\begin{center}\small\ttfamily a845ab5d0f4591294ef3c2d92feaeb5825a8d543.\end{center}
Other overview files were read on the live branch. This was not an exhaustive
repository audit, a build, or a verification of every theorem mentioned in
its documentation.

\wtag\ [At the write, the status file read at this commit is unchanged (same
blob, Appendix~\ref{sdf:app:provenance}). Two of the frontend tools used
below are already in the collection's report
\emph{canonical-diophantine-certificates}, and a non-claim of the report
\emph{probabilistic-quantum-and-continuous-computation} bears on the main
theorem; the \hyperref[sdf:sec:preface]{editorial preface} gives the
pointers.]

\subsection{Known antecedent and the precise novelty boundary}\label{sdf:sec:novelty}

Poonen proved that integer solvability can be preserved while replacing an
equation by one defining a smooth geometrically integral affine variety over
$\Q$. His Lemma~4.1 uses a positive multiple of $y^2-y$ added to a square,
with a suitable coefficient chosen effectively. He also obtained
undecidability under a smooth-projective-closure condition
\cite[Lemma 4.1 and Corollary 4.3]{poonen}.

Consequently, \emph{smoothness does not make Hilbert's tenth problem decidable}
is not a new claim. The proposed contribution here is the simultaneous,
coefficient-independent package: five additional variables, a quartic after
quadratic lifting, an exact natural-fibre bijection, smoothness over the entire
integer base, an integral Jacobian certificate of bounded degree, and explicit
dyadic and local points. A comprehensive priority search for this exact
combination was not completed. The construction is presented as a research
result with proofs, not as a certified world-first theorem.

\wtag\ [Variable count: Poonen's lemma adds one variable to a single
equation $f$ and gives degree $2\deg f$ and smoothness over~$\Q$; applied to
$f=\sum_if_i^2$ for a quadratic system it would give degree eight. The price
here of smoothness over all of~$\Z$, with an explicit certificate, degree
four and a natural-zero bijection, is four further variables. The
Hilbert's-tenth review~\cite{review13} also read Poonen's Lemma~4.1 and
Corollary~4.3 and found that they support this attribution and do not
contain the five-variable formula; it was not a complete priority search.]

\subsection{Reading map}

Sections~\ref{sdf:sec:main}--\ref{sdf:sec:integrality} prove the algebraic finalizer.
Sections~\ref{sdf:sec:local}--\ref{sdf:sec:real} prove the arithmetic and real-geometric
claims. Section~\ref{sdf:sec:guards} explains why three weighted guards are the
smallest successful choice in the specified weight-four architecture.
Sections~\ref{sdf:sec:counter}--\ref{sdf:sec:undecidable} connect the finalizer to
computation, carefully separating external time bounds from fixed-arity
unbounded representations. The implementation and its limits are recorded
in Section~\ref{sdf:sec:implementation}.

\section{Conventions and the quadratic interface}\label{sdf:sec:interface}

Throughout, $\N=\{0,1,2,\ldots\}$. A polynomial with integer coefficients is
evaluated in $\Z$ even when its variables range over $\N$. Negative
coefficients in a residual do not change the domain of its witnesses.
For a finite system $\source=(f_1,\ldots,f_m)$, put
\[
 \zeros{D}{\source}=
 \{a\in D^n:f_i(a)=0\text{ for every }i\},
 \qquad D\in\{\N,\Z,\Q,\R\}.
\]
Empty systems and zero residuals are allowed. Parameters are coefficients
unless explicitly declared to be additional variables.

\begin{definition}[Quadratic interface]
A quadratic certificate consists of named variables $x=(x_1,\ldots,x_n)$,
a finite list of polynomials $f_i\in\Z[x]$ of total degree at most two,
and a separately justified interpretation of its zeros over a stated domain.
The interpretation is not part of the algebraic data.
\end{definition}

\subsection{Functional arithmetic-circuit lifting}

\begin{lemma}[Quadratic lifting]\label{sdf:lem:lifting}
A finite system of integer polynomial equations can effectively be replaced
by a quadratic system whose projection onto the original variables is a
bijection of integer zero sets. Over $\N$, the same statement holds using
nonnegative arithmetic circuits for the positive and negative parts of each
polynomial.
\end{lemma}
\begin{proof}
Over $\Z$, evaluate each polynomial by a finite straight-line circuit. For an
addition, subtraction, or multiplication gate introduce a fresh variable $w$
and the corresponding equation
\[
 w-a-b=0,\qquad w-a+b=0,\qquad w-ab=0.
\]
Constants are imposed by $w-c=0$, and output variables are set to zero.
Every source tuple has exactly one gate assignment, by induction on the
acyclic circuit order. The residuals have degree at most two, and the outputs
vanish exactly when the original equations hold.

Over $\N$, first write $f=f^+-f^-$ with nonnegative integer coefficients.
Evaluate $f^+$ and $f^-$ using addition and multiplication only, and equate
the two output variables. All gate values are then natural, and the same
induction proves uniqueness. A circuit using arbitrary subtraction is not
silently treated as a natural circuit.
\end{proof}

This lemma preserves \emph{gate} uniqueness. It does not make the witnesses
of a pre-existing Diophantine representation unique. For example, a
representation of halting may have many auxiliary solutions before any
circuit gates are introduced. That multiplicity survives lifting.

\wtag\ [The natural version of Lemma~\ref{sdf:lem:lifting} is the
unique-extension quadratization of the collection's report
\emph{canonical-diophantine-certificates}, Lemma
\path{cdc:wf:lem:quadratization}, with the same proof: circuits for $f^+$
and $f^-$, one natural variable per gate, uniqueness by induction along the
circuit. The integer version adds subtraction gates to the same induction.
No novelty is claimed for either; the lemma is printed because
Theorem~\ref{sdf:thm:hardness} uses both versions.]

\subsection{Affine smoothness and degree conventions}

Our geometric object is the affine scheme
\[
 \variety=\Spec\bigl(\Z[x,t,u,z_1,z_2,z_3]/(F)\bigr).
\]
Smoothness over $\Z$ means that its structure morphism to $\Spec\Z$ is
smooth, not merely that its real locus has no singularities. A geometric
fibre is obtained by passing to an algebraic closure of a residue field of
$\Z$, including the generic characteristic-zero fibre.

We use the standard hypersurface criterion: if the ideal generated by $F$
and all its partial derivatives is the whole polynomial ring, the
derivative-open subsets of the hypersurface are standard smooth and cover
it \cite{stacks}. In particular, an integral identity expressing $1$ in that
ideal certifies relative smoothness in every characteristic.

Degree four will always mean degree in the \emph{witness variables of a
specialized equation}. When input parameters become coordinates of a single
master polynomial, their contribution to total degree must be counted again;
Section~\ref{sdf:sec:parameter-degree} does so explicitly.

\section{The five-variable finalizer}\label{sdf:sec:main}

Write each source residual as
\[
 f_i(x)=\sum_{|\alpha|\le2}c_{i,\alpha}x^\alpha.
\]
Its homogeneous quadratic lift is
\begin{equation}\label{sdf:eq:q}
 q_i(x,t)=\sum_{|\alpha|\le2}c_{i,\alpha}x^\alpha t^{2-|\alpha|}.
\end{equation}
The expression $t^2 f_i(x/t)$ is a useful mnemonic for $t\ne0$; the
polynomial definition \eqref{sdf:eq:q}, valid also at $t=0$, is authoritative.
Set
\begin{equation}\label{sdf:eq:ingredients}
 H(x,t)=\sum_{i=1}^m q_i(x,t)^2,\qquad
 R=tu-1,\qquad \guard(z)=2z^2-z.
\end{equation}
Put $\lambda=(\lambda_1,\lambda_2,\lambda_3)=(1,1,2)$. The finalizer is
\begin{equation}\boxed{\displaystyle
 F(x,t,u,z)=H(x,t)+(tu-1)^2+
             \guard(z_1)+\guard(z_2)+2\guard(z_3).}\label{sdf:eq:F}
\end{equation}
Here $z=(z_1,z_2,z_3)$, and all five new variables are distinct from the
source variables.

\begin{theorem}[Canonical smooth quartic finalizer]\label{sdf:thm:main}
For every quadratic interface in $n$ variables, the polynomial
\eqref{sdf:eq:F} has the following properties.
\begin{enumerate}[label=\textnormal{(\roman*)}]
\item Its coefficients are integers and its total degree is exactly four.
\item The map $a\mapsto(a,1,1,0,0,0)$ is a bijection from
$\zeros{\N}{\source}$ to $\zeros{\N}{F}$.
\item The map
\[
 (a,\varepsilon)\longmapsto
 (\varepsilon a,\varepsilon,\varepsilon,0,0,0)
\]
is a bijection from $\zeros{\Z}{\source}\times\{1,-1\}$ to
$\zeros{\Z}{F}$.
\item The affine scheme $\variety$ is smooth over $\Z$ of relative dimension
$n+4$, and every geometric fibre is integral.
\item An explicit identity $AF+\sum_v B_vF_v=1$ exists with
$\deg A\le3$ and $\deg B_v\le4$; its coefficients are given below without
an elimination or genericity search.
\item A point of $\variety(\Z[1/2])$ with nonnegative rational coordinates
can always be constructed. Also $\variety(\Z_p)\ne\varnothing$ for every
prime $p$, and $\variety(\Z/M\Z)\ne\varnothing$ for every $M\ge1$.
\item $\variety(\Q)$ is dense in $\variety(\R)$ and is Zariski dense in
$\variety_\Q$. The real locus has exactly two path components.
\end{enumerate}
\end{theorem}

The theorem is proved in the next five sections. Notice how different its
claims are: the first three are statements about an integer lattice, the
fourth and fifth about schemes, and the last two about rational, real, and
local points. None is being substituted for another.

\begin{proposition}[Size bounds]\label{sdf:prop:size}
If $f_i$ has $s_i$ nonzero monomials, the number of monomials of $F$ is at
most
\[
 9+\sum_{i=1}^m\binom{s_i+1}{2}.
\]
For the sum of absolute values of coefficients,
\[
 \lVert F\rVert_1\le \sum_i\lVert f_i\rVert_1^2+16.
\]
The construction of the expanded polynomial has polynomial bit complexity
in the expanded quadratic input size.
\end{proposition}
\begin{proof}
Homogenization preserves the number and coefficients of source monomials.
A square of an $s_i$-term polynomial contributes at most
$s_i(s_i+1)/2$ monomials. The unit square contributes at most three and
the three guards contribute six. Possible coincidences only lower the
count. The $\ell^1$ norm is submultiplicative; the unit square has norm four
and the three weighted guards together have norm twelve. The stated expansion uses a
quadratic number of input-term products with polynomial-length coefficients.
\end{proof}

These are expanded-polynomial bounds, not claims about the smallest arithmetic
circuit or the repository's universal operation-count records.

\section{Exact integer fibres and witness heights}\label{sdf:sec:fibres}

\begin{lemma}[A one-zero integer guard]\label{sdf:lem:guard}
For every $z\in\Z$, $\guard(z)\ge0$, with equality if and only if $z=0$.
For every $z\in\R$,
\begin{equation}\label{sdf:eq:guard-complete}
 \guard(z)=2\left(z-\frac14\right)^2-\frac18.
\end{equation}
\end{lemma}
\begin{proof}
For $z\ge1$, $z(2z-1)>0$; for $z\le-1$, both factors are negative, so
again their product is positive. The remaining integer is zero.
Completing the square gives \eqref{sdf:eq:guard-complete}.
\end{proof}

\begin{proof}[Proof of Theorem~\ref{sdf:thm:main}(i)--(iii)]
Each $q_i$ is homogeneous of degree two, so $H$ has degree at most four.
The monomial $t^2u^2$ has coefficient one in $F$: no term of $H$ involves
$u$, and the guards involve neither $t$ nor $u$. Thus the degree is exactly
four, including for the empty or identically zero source system.

On an integer tuple, every summand in \eqref{sdf:eq:F} is nonnegative.
Consequently $F=0$ forces
\[
 q_i(x,t)=0\ (1\le i\le m),\qquad tu=1,\qquad z_1=z_2=z_3=0.
\]
The units of $\Z$ are $\pm1$, so $t=u=\varepsilon$ for one
$\varepsilon\in\{1,-1\}$. Put $a=\varepsilon x$. Homogeneity gives
\[
 q_i(\varepsilon a,\varepsilon)=\varepsilon^2 f_i(a)=f_i(a).
\]
Thus $a$ is a source zero. Conversely, every displayed lift of a source
zero makes every summand vanish. The inverse map is
$(x,t,u,z)\mapsto(tx,t)$, proving bijectivity and not just existence.

If all coordinates are natural, $tu=1$ forces $t=u=1$. The resulting
inverse map is simply projection onto $x$. No sign choice remains.
\end{proof}

\begin{corollary}[Exact counting and height transfer]\label{sdf:cor:height}
Let $B\ge1$ be an integer, and count zeros whose coordinates have absolute
value at most $B$. The number of target natural zeros is exactly the number
of source natural zeros; the number of target integer zeros is twice the
number of source integer zeros. For each individual lift,
\[
 \height(\varepsilon a,\varepsilon,\varepsilon,0^3)
       =\max\{1,\height(a)\},\qquad
 \height(a)=\max_i|a_i|.
\]
The same assertions hold fibrewise after any external parameters are fixed.
\end{corollary}
\begin{proof}
The bijections just proved add only $\pm1$ and zero and change source
coordinates only by a common sign. They therefore preserve every box of
radius $B\ge1$ in the asserted manner.
\end{proof}

\begin{remark}[What the extra coordinates do not encode]
The three $z$ coordinates are zero at every integer solution. They are not
additional computational choices. The $t,u$ coordinates impose a unit
condition, and over $\N$ that unit is uniquely one. Their geometric role
becomes visible only away from the integer lattice, where the guards may
be negative and $t$ may be a nonintegral unit.
\end{remark}

\section{Smoothness in every characteristic}\label{sdf:sec:smooth}

There are two complementary proofs. The first explains the choice of
constants through critical values. The second gives an integral identity
suitable for independent symbolic checking and eventual formalization.

\subsection{The critical-value mechanism}

Put $Q=H+R^2$, so $F=Q+\sum_j\lambda_j\guard(z_j)$.

\begin{lemma}[Only two possible critical values]\label{sdf:lem:critical}
Over any field of characteristic different from two, every critical point
of $Q$ has value either zero or one.
\end{lemma}
\begin{proof}
Homogeneity of $H$ gives the polynomial identity
\begin{equation}\label{sdf:eq:Euler}
 \sum_{i=1}^n x_iQ_{x_i}+tQ_t-uQ_u=4H.
\end{equation}
The contributions of $R^2$ cancel between the last two terms. At a critical
point the left side is zero, hence $H=0$. Moreover $Q_u=2tR$.
If $t\ne0$, then $R=0$ and $Q=0$. If $t=0$, then $R=-1$ and $Q=1$.
\end{proof}

A singular point of $F=0$ in odd characteristic would satisfy
$4z_j-1=0$ for all three guards. Their weighted values sum to $-4/8$, so
$Q=1/2$. That contradicts Lemma~\ref{sdf:lem:critical}. In characteristic two,
$F_{z_1}=4z_1-1=1$, which rules out singularities directly.

This calculation is coefficient independent. In particular, no assumption
about the source zero set, its dimension, or the regularity of its defining
system is used.

\subsection{An integral unit-ideal certificate}

For derivatives of the full polynomial $F$, define
\begin{equation}\label{sdf:eq:EDJ}
 \begin{split}
 E&=\sum_{i=1}^n x_iF_{x_i}+tF_t-uF_u=4H,\\
 D_j&=F_{z_j}=\lambda_j(4z_j-1),\\
 J&=8F-2E-\sum_{j=1}^3(4z_j-1)D_j=8R^2-4.
 \end{split}
\end{equation}
The rightmost equalities are identities over $\Z$, not reductions on the
zero set.

\begin{proposition}[Universal Jacobian identity]\label{sdf:prop:jacobian}
In $\Z[x,t,u,z]$,
\begin{equation}\label{sdf:eq:four}
 4u(2R-1)F_u-(2R+1)J=4
\end{equation}
and therefore
\begin{equation}\boxed{\displaystyle
 z_1\bigl[4u(2R-1)F_u-(2R+1)J\bigr]-D_1=1.}\label{sdf:eq:unit}
\end{equation}
In particular, $1$ belongs to the ideal generated by $F$ and its partial
derivatives.
\end{proposition}
\begin{proof}
Substitute $F_u=2tR$ and $J=8R^2-4$ into the left side of
\eqref{sdf:eq:four}, and use $tu=R+1$. It becomes
\[
 8R(R+1)(2R-1)-(2R+1)(8R^2-4)=4.
\]
Multiplying by $z_1$ and subtracting $D_1=4z_1-1$ gives
\eqref{sdf:eq:unit}. By \eqref{sdf:eq:EDJ}, both $E$ and $J$ are explicit
combinations of the equation and its derivatives. Every product $(4z_j-1)D_j$ is a derivative times a polynomial, so it
also belongs to the same ideal.
\end{proof}

For direct certificate exchange, no nested abbreviations are necessary.
Let $F_v=\partial F/\partial v$, and set
\begin{equation}\label{sdf:eq:coefficients}
\begin{aligned}
 A&=-8z_1(2R+1),\\
 B_{x_i}&=2z_1(2R+1)x_i,&
 B_t&=2z_1(2R+1)t,\\
 B_u&=2z_1u(2R-3),&
 B_{z_j}&=z_1(2R+1)(4z_j-1)-\delta_{j1}.
\end{aligned}
\end{equation}
Expansion of \eqref{sdf:eq:unit} gives
\begin{equation}\boxed{\displaystyle AF+\sum_v B_vF_v=1.}\label{sdf:eq:bezout}
\end{equation}
Here $\deg A\le3$ and $\deg B_v\le4$, so every summand in this certificate
has degree at most seven. The degree bound is independent of the number of
source equations and the size of their coefficients.

\begin{proof}[Proof of smoothness in Theorem~\ref{sdf:thm:main}]
The identity \eqref{sdf:eq:bezout} persists under every base change. On the
hypersurface, the partial derivatives generate the unit ideal. The open
subsets where one partial is invertible form a standard-smooth cover
\cite{stacks}. Thus $\variety\to\Spec\Z$ is smooth. It is locally one
smooth equation in $n+5$ variables, so its relative dimension is $n+4$.
This argument does not infer relative smoothness merely from a picture of
the real locus or from a finite list of prime checks.
\end{proof}

\begin{remark}[Universal coefficient families]
The same identity works with coefficients in any commutative base ring.
In particular, letting the source coefficients vary gives a smooth relative
family over the coefficient-parameter scheme. This does not assert that
there is a polynomially varying dyadic point over that entire parameter
space; the point construction below is performed after integer
specialization.
\end{remark}

\section{Geometric integrality and the special fibre at two}\label{sdf:sec:integrality}

\begin{proposition}\label{sdf:prop:integral}
Every geometric fibre of $\variety\to\Spec\Z$ is integral.
In characteristic two the fibre is isomorphic to affine $(n+4)$-space.
\end{proposition}
\begin{proof}
Let $k$ be an algebraically closed field of odd characteristic or
characteristic zero. Regard $F$ as a polynomial in $z_1,z_2,z_3$ over
$K=k(x,t,u)$. Its highest homogeneous part in the $z$ variables is
\[
 2z_1^2+2z_2^2+4z_3^2,
\]
a quadratic form of rank three. A nontrivial factorization of this
degree-two polynomial over $K$ would have two affine-linear factors in
$z$. Their leading product has rank at most two, a contradiction.
Hence $F$ is irreducible over $K$.

As a polynomial in $z$ over $k[x,t,u]$, $F$ is primitive: the coefficient
of $z_1^2$ is the nonzero constant $2$. Gauss's lemma now gives
irreducibility in $k[x,t,u,z]$, a unique factorization domain. The
principal ideal $(F)$ is prime, so the fibre is integral.

In characteristic two, every $2z_j^2$ disappears and the coefficient of
$z_1$ is $-1=1$. Solving uniquely for $z_1$ identifies the coordinate ring
with a polynomial ring in the remaining $n+4$ variables. This proves both
integrality and the stated affine-space description.
\end{proof}

There is no hidden component carrying the actual computation while another
component supplies easy rational points. The entire geometric fibre is one
integral variety. The integer and rational behaviours coexist on that single
variety.

The phrase ``good reduction'' in this article refers to the displayed
\emph{affine} smooth model. It does not say that the naive projective closure
is smooth, or that there is a smooth proper model over $\Z$.

\section{Dyadic points and the complete absence of local obstructions}\label{sdf:sec:local}

\subsection{An explicit point after inverting only two}

Let
\begin{equation}\label{sdf:eq:c}
 c=\sum_{i=1}^m f_i(0)^2\in\N.
\end{equation}
Choose $K=2^k$ with $K^4>2c$, and set $N=K^4-2c>0$.
Lagrange's four-square theorem supplies nonnegative integers
$a_1,\ldots,a_4$ such that
\begin{equation}\label{sdf:eq:four-square}
 N=a_1^2+a_2^2+a_3^2+a_4^2.
\end{equation}
We use this classical theorem only for a finite integer, not as a statement
about representing computations; it is also the standard bridge from
natural to integer witnesses in Diophantine undecidability arguments
\cite[Section 5]{poonen}.

\begin{lemma}[An elementary universal weighted form]\label{sdf:lem:1122}
Every nonnegative integer is of the form
\[
 A_1^2+A_2^2+2A_3^2+2B^2,
 \qquad A_1,A_2,A_3,B\in\N.
\]
The same form represents every nonnegative rational number over $\Q$.
\end{lemma}
\begin{proof}
Start from four ordinary squares. Among four integers, two have the same
parity. Replace the squares of that pair, say $a,b$, by
\[
 a^2+b^2=2\left(\frac{a+b}{2}\right)^2+
        2\left(\frac{a-b}{2}\right)^2.
\]
The other two squares remain unchanged, and an absolute value makes the last
coordinate nonnegative. For $A/B\ge0$, represent the integer $AB$ and
divide all four coordinates by $B$.
\end{proof}

\wtag\ [The form $x^2+y^2+2z^2+2w^2$ is one of the universal positive
diagonal quaternary forms listed by Ramanujan~\cite{ramanujan}; the
elementary parity proof above derives its universality from Lagrange's
theorem.]

\begin{theorem}[Constructive dyadic solubility]\label{sdf:thm:dyadic}
Choose a representation
\[
 K^4-2c=A_1^2+A_2^2+2A_3^2+2B^2.
\]
Then the coordinates
\begin{equation}\label{sdf:eq:dyadic}
 x=0,\qquad t=\frac1K,\qquad u=K+\frac BK,\qquad
 z_j=\frac14+\frac{A_j}{2K^2}\quad(1\le j\le3)
\end{equation}
define a point of $\variety(\Z[1/2])$, all of whose coordinates are
nonnegative. If $K$ is the smallest admissible power of two, the rational
coordinate height is $O(\sqrt{c+1})$, and its bit length is $O(\log(c+2))$.
\end{theorem}
\begin{proof}
At these coordinates $R=B/K^2$ and $H=c/K^4$. Completing the weighted
squares gives
\[
 \sum_{j=1}^3\lambda_j\guard(z_j)
 =\frac{A_1^2+A_2^2+2A_3^2}{2K^4}-\frac12.
\]
Adding $R^2$ and $H$ yields zero by the displayed representation of
$K^4-2c$. All denominators are powers of two. The coordinates satisfy
$1/4\le z_j\le3/4$ and $K\le u\le2K$.

For a reduced rational $a/b$ with $b>0$, use coordinate height
$\max\{|a|,b\}$. All coordinates in \eqref{sdf:eq:dyadic} have height at most
$4K^2$. If $c\ge1$, minimality gives $(K/2)^4\le2c$, hence $K^4\le32c$.
The case $c=0$ has $K=1$. These inequalities give the claimed bounds.
\end{proof}

The small height bound is not a deterministic polynomial-time claim for
finding a four-square decomposition in the bit model. The included reference
routine uses a finite pair-table search, and explicitly limits the size of
that search. A supplied decomposition can instead be checked immediately
by \eqref{sdf:eq:four-square}.

\subsection{Integral points at all primes}

\begin{theorem}[Every integral local test succeeds]\label{sdf:thm:local}
For every prime $p$, $\variety(\Z_p)$ is nonempty. Consequently the equation
has a solution modulo every positive integer.
\end{theorem}
\begin{proof}
For an odd prime $p$, the dyadic point \eqref{sdf:eq:dyadic} is already integral
in $\Z_p$, because its denominators are powers of two.

For $p=2$, fix $x=0$, $t=u=1$, and $z_2=z_3=0$. The remaining equation
is
\begin{equation}\label{sdf:eq:2adic}
 h(z_1):=c+2z_1^2-z_1=0.
\end{equation}
Modulo two it has the root $z_1\equiv c$, and $h'(z_1)=4z_1-1$ is odd.
Hensel lifting therefore gives a root in $\Z_2$. To avoid treating the
lifting as a black box, suppose $h(z_s)\equiv0\pmod{2^s}$ and choose
$b\in\{0,1\}$ so that
\[
 h(z_s)/2^s+b\equiv0\pmod2.
\]
Then $z_{s+1}=z_s+b2^s$ satisfies
$h(z_{s+1})\equiv0\pmod{2^{s+1}}$. The compatible residues converge to the
required root.

For $M=\prod p^{e_p}$, reduce one local point modulo each $p^{e_p}$ and
apply the Chinese remainder theorem coordinate by coordinate. The resulting
point solves $F=0$ modulo $M$.
\end{proof}

Thus rejecting source instances cannot be rejected by finding a bad prime,
a bad prime power, or any congruence modulus. They even have compatible
integral points at every prime. This is a failure of a prospective
\emph{integral} local-to-global principle, not a failure of the rational
Hasse principle: rational points are already available globally.

\subsection{A five-variable counterexample requiring no machine encoding}

\begin{example}[An everywhere soluble equation with no integer point]\label{sdf:ex:one}
Take the inconsistent source equation $1=0$, with no source variables.
Then
\begin{equation}\label{sdf:eq:inconsistent}
 t^4+(tu-1)^2+\guard(z_1)+\guard(z_2)+2\guard(z_3)=0
\end{equation}
has no integer solution. Indeed, vanishing forces $tu=1$ and $t^4=0$,
which are incompatible. Nevertheless it has the dyadic point
\[
 \left(t,u,z_1,z_2,z_3\right)
 =\left(\frac12,3,\frac14,\frac12,\frac38\right).
\]
Here $K=2$ and $14=0^2+2^2+2\cdot1^2+2\cdot2^2$.
The Jacobian identity proves smoothness over $\Z$, not just at this point;
Theorem~\ref{sdf:thm:local} supplies all integral local points.
\end{example}

This example is elementary, but its role in the general reduction is
important. The same harmless arithmetic geometry occurs for instances whose
integer solvability encodes arbitrary halting behaviour. The phenomenon is
not restricted to isolated inconsistent equations that an algorithm could
recognize syntactically.

\section{Rational density and the real locus}\label{sdf:sec:real}

\subsection{Rational points in the quadric fibres}

At every real point, completing the weighted guard squares gives
\[
 H+(tu-1)^2\le\frac12,
 \qquad tu\ge1-\frac1{\sqrt2}>0.
\]
In particular $t$ is nonzero. We may therefore project onto the base
$b=(x,t)$ and recover $u=(R+1)/t$ from $R$. Put
\begin{equation}\label{sdf:eq:rho}
 \rho(b)=\frac12-H(x,t),\qquad y_j=z_j-\frac14.
\end{equation}
The remaining equation is the ellipsoid
\begin{equation}\label{sdf:eq:sphere}
 R^2+2y_1^2+2y_2^2+4y_3^2=\rho(b).
\end{equation}
For rational $b$ with $\rho(b)>0$, this quadric has a rational point.
Indeed, after multiplication by two, its left side becomes
\[
 (2y_1)^2+(2y_2)^2+2(2y_3)^2+2R^2,
\]
and Lemma~\ref{sdf:lem:1122} represents the positive rational $2\rho(b)$.

A nonsingular ellipsoid with rational coefficients and a rational point has
dense rational points in its real locus. A line through the rational point
with rational direction meets the quadric in a second rational point:
substitution gives a quadratic equation with one known root. The resulting
rational parametrization covers the quadric except for the usual
lower-dimensional exceptional set, and rational parameters are dense in
real parameters. Equivalently, use stereographic projection after fixing
a rational point and its rational tangent hyperplane.

\begin{theorem}[Real and Zariski density]\label{sdf:thm:density}
The set $\variety(\Q)$ is dense in $\variety(\R)$ in the Euclidean topology
and is Zariski dense in $\variety_\Q$.
\end{theorem}
\begin{proof}
Consider a real point with $\rho(b)>0$. Approximate $b=(x,t)$ by rational
base points $b_\nu$ with nonzero $t_\nu$ and positive $\rho$. Radially
rescale $(R,y_1,y_2,y_3)$ to get nearby real points on the nearby ellipsoids.
On each such ellipsoid choose a rational point arbitrarily close to that
real point. Recovering $u=(R+1)/t$ is continuous because $t\ne0$. This
produces rational points tending to the original point.

At a boundary point $\rho=0$, all four coordinates $(R,y)$ vanish and
$H=1/2$. Euler's identity
$\sum x_iH_{x_i}+tH_t=4H$ shows that $\nabla H$ is not zero there.
Thus every neighbourhood contains base points with $\rho>0$. The ellipsoid
radii tend to zero on approach to the boundary, so the same approximation
also reaches the boundary point.

There is a nonempty real open set with $t\ne0$ and $\rho>0$: take $x=0$
and positive $t$ sufficiently small. The smooth integral variety therefore
has real manifold points of its full dimension $n+4$. A polynomial function
vanishing on all its rational points vanishes, by Euclidean density, on this
real open manifold. In smooth local algebraic coordinates a regular
function vanishing on a real open set vanishes identically on the
irreducible variety. Hence the rational points are Zariski dense.
\end{proof}

The last local fact can also be seen by complexifying a smooth coordinate
chart: a holomorphic function that vanishes on an open subset of its real
coordinate slice has identically zero Taylor series. Geometric integrality
then propagates the identity to the whole variety. No density assertion is
made for \emph{dyadic} points; ordinary rational directions, with arbitrary
denominators, were used in the quadric argument.

\subsection{Exactly two path components}

\begin{theorem}\label{sdf:thm:components}
The real locus has exactly two path components, distinguished by the sign
of $t$, or equivalently the sign of $u$.
\end{theorem}
\begin{proof}
On $\variety(\R)$, the sum of the guards is at least $-1/2$. Hence
\[
 H+(tu-1)^2\le\frac12,
 \qquad tu\ge1-\frac1{\sqrt2}>0.
\]
Thus $t,u$ are nonzero and have the same sign. Both signs occur, for example
by changing $(x,t,u)$ to $(-x,-t,-u)$ in the point
\eqref{sdf:eq:dyadic}. No continuous path can join the two signs.

To show that the positive part is path connected, use coordinates
\[
 y=x/t,\qquad r=tu-1,\qquad w_j=\sqrt{2\lambda_j}(z_j-1/4),\qquad t>0.
\]
If $S(y)=\sum_i f_i(y)^2$, the equation becomes
\begin{equation}\label{sdf:eq:real-model}
 t^4S(y)+r^2+\sum_{j=1}^3w_j^2=\frac12.
\end{equation}
Conversely, every solution of \eqref{sdf:eq:real-model} with $t>0$ gives an
original point: $u=(1+r)/t>0$ since $|r|\le1/\sqrt2$.

For fixed $y$, decreasing $t$ makes the squared radius
$1/2-t^4S(y)$ nondecreasing. A point can therefore be moved to arbitrarily
small positive $t$ by radially rescaling its vector $(r,w)$. If that vector
is initially zero, choose any direction when the radius first becomes
positive; the resulting path is still continuous at the initial point.
For a positive radius its direction lies in the path-connected sphere
$S^3$, so the direction may be changed to any chosen common direction.

Now join $y$ to $0$ by a line segment. The continuous function $S$ is bounded
on this compact segment. Choose $t=\epsilon>0$ small enough that
$\epsilon^4S<1/2$ throughout it. Keeping the common direction and varying
the radius continuously gives a path along the whole segment. At $y=0$,
any two admissible positive values of $t$ can be joined, again by radial
rescaling. Thus all positive-$t$ points lie in one path component. The sign
symmetry proves the same for negative $t$.
\end{proof}

Only the number of components is uniform. For example, if $H=0$, each
component has an $S^3$ factor in the real description. If the only residual
is the constant $1$, the corresponding radius closes at a finite positive
$t$. We do not assert a single source-independent diffeomorphism type.

\section{Why three weighted guards: a sharp restricted obstruction}\label{sdf:sec:guards}

The five-variable result is not obtained merely by deleting one guard from
a sum of four identical ones. The weights and the use of the unit residual
in the rational-point construction are essential.

\subsection{Classification by total guard weight}

Consider the more general architecture
\begin{equation}\label{sdf:eq:Fweights}
 F_{\mathbf w}=H+(tu-1)^2+\sum_{j=1}^k w_j(2z_j^2-z_j),
 \qquad w_j\in\Z_{>0},\quad w_1=1,
\end{equation}
and let $W=\sum_jw_j$. The assumption $w_1=1$ gives an especially simple
uniform certificate at the prime two. Define
\[
 E=4H,\qquad D_j=w_j(4z_j-1),\qquad
 J_W=8F_{\mathbf w}-2E-\sum_j(4z_j-1)D_j=8R^2-W.
\]
A general elimination identity is
\begin{equation}\label{sdf:eq:general-guard}
 W(8-W)=4u(8R-W)(F_{\mathbf w})_u-(8R+8-W)J_W.
\end{equation}
It follows by substituting $(F_{\mathbf w})_u=2tR$ and $tu=R+1$.

\begin{theorem}[Universal smoothness in the weighted architecture]\label{sdf:thm:weights}
Under $w_1=1$, the hypersurface \eqref{sdf:eq:Fweights} is smooth over $\Z$
for \emph{every} quadratic source system if and only if
\[
 W\in\{4,16\}.
\]
\end{theorem}
\begin{proof}
In the quotient by the equation and all its derivatives, $D_1=0$ gives
$4z_1=1$, so two is invertible. Identity \eqref{sdf:eq:general-guard} says that
$W(8-W)$ is zero in that quotient. For $W=4$ or $W=16$ it is a nonzero
signed power of two. The quotient is therefore the zero ring, proving that
the Jacobian ideal is the unit ideal and the hypersurface is smooth.

For necessity, suppose an odd prime $p$ divides $8-W$. At the point
$x=t=u=0$ and $z_j=1/4$ over $\overline{\Fp}$, the equation and all its
partial derivatives vanish. This works for every source, because $H$ is
homogeneous of degree four and has zero gradient at the origin. The case
$W=8$ gives the same counterexample for every odd prime.

If instead an odd prime $p$ divides $W$, choose the zero source $H=0$ and
take $t=u=1$, $x=0$, and $z_j=1/4$. This is again a singular point. Hence
universal smoothness requires that both $W$ and $|8-W|$ be nonzero powers
of two. Write $W=2^a$. If $a<3$, the odd factor of $8-W$ is
$2^{3-a}-1$, which equals one only for $a=2$. If $a>3$, the odd factor of
$W-8$ is $2^{a-3}-1$, which equals one only for $a=4$. The case $a=3$
was excluded. Thus $W=4$ or $16$.
\end{proof}

For identical guards $w_j=1$, this says that exactly $k=4$ and $k=16$
give universal smoothness. In particular, simply using one, two, or three
identical guards always introduces a bad odd prime: respectively seven,
three, or five at the origin. The weighting $(1,1,2)$ keeps total weight
four while using only three variables.

\subsection{Why the still smaller weighting \texorpdfstring{$(1,3)$}{(1,3)} is insufficient}

Two guards with weights $(1,3)$ also have total weight four, and therefore
satisfy the smoothness theorem. They do \emph{not}, however, give a dyadic
point for every source. This is not merely a failure of our chosen
point-construction algorithm.

\begin{lemma}[A dyadic ternary-form obstruction]\label{sdf:lem:dyadic-obstruction}
If an integer $N$ satisfies $N\equiv10\pmod{16}$, the equation
\[
 X^2+3Y^2+2Z^2=N
\]
has no solution with $X,Y,Z\in\Z[1/2]$.
\end{lemma}
\begin{proof}
First, such a solution would have integral coordinates. Otherwise multiply
by a smallest common denominator $2^d$, $d\ge1$, to obtain integers
$X_0,Y_0,Z_0$, not all even, with
\[
 X_0^2+3Y_0^2+2Z_0^2=2^{2d}N\equiv0\pmod8.
\]
Reduction modulo four forces $Z_0$ even and $X_0,Y_0$ of the same parity.
Primitivity then makes $X_0,Y_0$ odd, but the left side is four modulo
eight, a contradiction.

For an integral solution, the even value of $N$ forces $X,Y$ to have the
same parity. If both are odd, the left side is four or six modulo eight,
not $N\equiv2\pmod8$. If both are even, write $X=2a$, $Y=2b$.
Reduction modulo four makes $Z$ odd, so $2Z^2\equiv2\pmod{16}$.
The equation would imply $a^2+3b^2\equiv2\pmod4$, which is impossible.
\end{proof}

\begin{proposition}[Failure of the two-guard dyadic guarantee]\label{sdf:prop:two-guard}
For the source consisting of three constant residuals equal to one,
\begin{equation}\label{sdf:eq:bad-weighted}
 3t^4+(tu-1)^2+\guard(z)+3\guard(w)=0
\end{equation}
has no point over $\Z[1/2]$, although its affine model is smooth over $\Z$.
\end{proposition}
\begin{proof}
A hypothetical dyadic point is a real point. Completing squares gives
$3t^4\le1/2$, while the same estimate as before gives $tu>0$.
Thus $0<|t|<1$. Write $t=a/2^k$ with $a$ odd and $k\ge1$.
With $R=tu-1$, the equation becomes
\[
 (2z-1/2)^2+3(2w-1/2)^2+2R^2=1-6t^4.
\]
Multiplying all three coordinates on the left by $2^{2k}$ gives a dyadic
representation by $X^2+3Y^2+2Z^2$ of the integer
\[
 N=2^{4k}-6a^4\equiv10\pmod{16}.
\]
This contradicts Lemma~\ref{sdf:lem:dyadic-obstruction}. Smoothness follows
from total weight four and first weight one.
\end{proof}

\begin{corollary}[Five auxiliaries are optimal in the weight-four family]\label{sdf:cor:five}
Among finalizers of the form
\[
 H+(tu-1)^2+\sum_jw_j\guard(z_j),\qquad
 w_j>0,\quad \sum_jw_j=4,
\]
which must be smooth over $\Z$ and have a dyadic point for every source,
at least three guard variables are necessary. The choice $(1,1,2)$ is
sufficient, so five total auxiliary variables are optimal within this family.
\end{corollary}
\noindent\wtag\ [In Corollary~\ref{sdf:cor:five}, $w_j\in\Z_{>0}$, as in
\eqref{sdf:eq:Fweights}; the display ``$w_j>0$'' is the source's. The
normalization $w_1=1$ of Theorem~\ref{sdf:thm:weights} is not assumed: the
all-even weightings $(4)$ and $(2,2)$ are excluded by the all-even argument
of the proof below, and $(1,3)$ and $(1,1,2)$ have a weight one after
reordering.]
\begin{proof}
With one guard the only weight is four. With two guards the weights are,
up to order, $(2,2)$ or $(1,3)$. All-even weights give a zero gradient
modulo two for the zero source $H=0$, and a nonempty singular fibre.
The remaining pair $(1,3)$ is excluded by Proposition~\ref{sdf:prop:two-guard}.
Theorem~\ref{sdf:thm:main} proves sufficiency of three guards.
\end{proof}

This is a restricted optimality theorem. It does not rule out a different
unit gadget, a different nonnegative polynomial guard, a different total
weight, or a construction with fewer variables that avoids this architecture
entirely.

\section{A concrete Turing-complete frontend}\label{sdf:sec:counter}

The finalizer does not depend on a particular language. Nevertheless, an
executable frontend makes the semantics and the witness-count claims
concrete. We use finite deterministic counter programs with natural
registers and instructions of the following forms:
\[
 \operatorname{inc}(j,q),\qquad
 \operatorname{dec}(j,q_+,q_0),\qquad
 \operatorname{halt}.
\]
An increment adds one to counter $j$ and jumps to $q$. A decrement
instruction decrements a positive counter and jumps to $q_+$, or leaves a
zero counter unchanged and jumps to $q_0$. There is one halt state, and it
has an absorbing self-loop for padding.

\subsection{Why this instruction set can support universal computation}

For completeness, three counters suffice for a direct simulation of a
Turing machine. Store the tape to the left and to the right of the current
head in two stacks, keep the current head symbol in finite control, and
use a third counter as scratch space. Encode a stack in a fixed base $B$
using nonzero digits for tape symbols, with the empty stack encoded by
zero. The outermost stack symbol is the least significant digit.

To push a digit $d$, replace $X$ by $BX+d$. Repeatedly decrement $X$ while
adding $B$ units to a scratch counter; then transfer the scratch counter
back to $X$ and add $d$. The fixed additions by $B$ or $d$ expand into
finite sequences of increment instructions.

To pop, repeatedly decrement $X$, maintaining the remainder modulo $B$ in
finite control and incrementing the scratch counter after each complete
block of $B$ decrements. At the end, the scratch counter is the quotient
and the control state holds the remainder. Transfer the quotient back to
$X$. An empty stack returns the blank symbol. Scratch space is zero at the
end of either operation.

A tape move pushes the updated head symbol onto one stack and pops from
the other. Thus each transition expands into a finite counter program,
with the unused stack preserved. This supplies an explicit simulation
argument for universality of the program class. The small worked example
below uses only one counter and is not itself a universal machine.
Formalized reductions through Minsky machines and FRACTRAN provide an
independent established context for this substrate choice \cite{coq}.
\wtag\ [Minsky showed that two counters already suffice for universal
computation~\cite{minsky}; the three-counter simulation above is a direct
argument, not a sharper count.]

\subsection{A quadratic certificate for an external horizon}

Fix a program with $\ell$ states, $r$ counters, and a chosen horizon $T$.
Expand instructions into edges $e$: increment edges, positive-decrement
edges, zero-test edges, and the absorbing halt edge. Let $E$ be their
number, $p$ the number of positive-decrement edges, and $g$ the number of
all guarded edges (positive and zero). Write $s(e),t(e)$ for source and
target states, and $\delta_{e,j}\in\{-1,0,1\}$ for counter updates.

Introduce natural variables
\[
 q_{s,i}\quad(0\le s\le T,\ 0\le i<\ell),\qquad
 c_{s,j}\quad(0\le s\le T,\ 0\le j<r),
\]
edge selectors $a_{s,e}$ for $s<T$, and one slack $d_{s,e}$ for every
positive-decrement edge at every time $s<T$. The equations at time $s$ are
\begin{align}
 \sum_e a_{s,e}&=1,\label{sdf:eq:oneedge}\\
 q_{s,i}&=\sum_{e:s(e)=i}a_{s,e},&
 q_{s+1,i}&=\sum_{e:t(e)=i}a_{s,e},\label{sdf:eq:states}\\
 c_{s+1,j}&=c_{s,j}+\sum_e\delta_{e,j}a_{s,e}.\label{sdf:eq:update}
\end{align}
For an edge testing counter $j$, impose
\begin{align}
 a_{s,e}c_{s,j}&=0&&\text{on a zero edge},\label{sdf:eq:zeroedge}\\
 a_{s,e}c_{s,j}-a_{s,e}-d_{s,e}&=0&&\text{on a positive edge}.\label{sdf:eq:posedge}
\end{align}
The initial one-hot state and initial counters are fixed by linear
equations. Finally impose $q_{T,h}=1$, where $h$ is the halt state.
Every displayed residual is of degree at most two.

\begin{theorem}[Canonical bounded counter certificates]\label{sdf:thm:counter}
The natural zero set of this quadratic system is empty if the program has
not halted by time $T$, and is a singleton if it has. The same statement
holds for its smooth quartic finalization, with five additional variables.
The quadratic interface has
\begin{align}
 n_T&=(\ell+r)(T+1)+(E+p)T,\label{sdf:eq:nT}\\
 m_T&=(\ell+r)+T(1+2\ell+r+g)+1.\label{sdf:eq:mT}
\end{align}
\end{theorem}
\begin{proof}
Because the selectors are natural and sum to one, exactly one selector is
one and all the others are zero. The state equations identify its source
and target. No separate Boolean equation is needed. On a selected zero
edge, \eqref{sdf:eq:zeroedge} forces the counter to be zero. On a selected
positive edge, \eqref{sdf:eq:posedge} says $c_{s,j}=1+d_{s,e}$, so the
counter is positive. On an inactive positive edge, that same equation
forces its slack to be zero, rather than leaving an arbitrary witness.

Starting from the fixed initial state, exactly one operational edge is
enabled. Induction on $s$ therefore determines every selector, every
counter, and every state variable uniquely. Active positive slacks equal
the tested counter minus one; inactive slacks are zero. Conversely the
actual padded run provides precisely this assignment and satisfies all
transition equations. The endpoint equation holds exactly when the run
has reached the absorbing halt state by time $T$.

There are $(\ell+r)(T+1)$ state and counter variables, $ET$ selectors,
and $pT$ slacks. The initial rows number $\ell+r$. Each step has one
selector-sum row, $2\ell$ state rows, $r$ update rows, and $g$ guard
rows. There is one final endpoint row. This gives
\eqref{sdf:eq:nT}--\eqref{sdf:eq:mT}, and Theorem~\ref{sdf:thm:main} gives the
last assertion.
\end{proof}

The proof depends on natural coordinates. Over unrestricted integers a sum
of selectors equal to one does not make them one-hot. The integer
zero-set theorem for the finalizer still holds, but it refers to the
integer zeros of the source equations, which need not describe legal
counter executions. A separate domain encoding is required before using
this frontend as an integer execution certificate.

\wtag\ [Theorem~\ref{sdf:thm:counter} is a second presentation of the
counter-program adapter of \emph{canonical-diophantine-certificates}: its
one-hot control equations, the zero guard \eqref{sdf:eq:zeroedge}, which is
that report's equation \path{cdc:eq:zero-guard}, the externally fixed
horizon of its Section \path{cdc:pt:sec:boundary} (arity growing with
$T$), and the halt self-loop caveat of its Section
\path{cdc:ct:sec:substrates}. New here are only the exact counts
\eqref{sdf:eq:nT}--\eqref{sdf:eq:mT} and the positive-edge slack equation
\eqref{sdf:eq:posedge}, with natural selectors summing to one in place of
Boolean rows. The worked example below is printed as an illustration of the
finalizer.]

\subsection{Worked example: a countdown program}

Take two states: state zero executes
$\operatorname{dec}(0,0,1)$, and state one halts. The three edges are
$D$ (decrement and loop), $Z$ (zero-test and enter halt), and $H$ (halt
self-loop). For input counter value two, the run is
\[
 (0,2)\xrightarrow{D}(0,1)\xrightarrow{D}(0,0)
       \xrightarrow{Z}(1,0)\xrightarrow{H}(1,0)\longrightarrow\cdots.
\]
Thus horizon two rejects; horizon three accepts; larger horizons add only
forced halt padding. Here
\[
 \ell=2,\quad r=1,\quad E=3,\quad p=1,\quad g=2,
 \qquad n_T=7T+3,\quad m_T=8T+4.
\]
For $T=3$ the source has 24 variables and 28 residuals. The exported
smooth equation has 29 variables, exact degree four, and 97 nonzero
monomials. Its natural zero set is a singleton.

The exported rational point deliberately need not coincide with the
computational witness. The source constant statistic is $c=9$; the
reference constructor takes $K=4$ and uses
\[
 238=0^2+2^2+3^2+15^2=3^2+15^2+2\cdot1^2+2\cdot1^2.
\]
All 24 source coordinates are set to zero in that rational point, with
\[
 t=\frac14,\quad u=\frac{17}{4},\quad
 (z_1,z_2,z_3)=\left(\frac{11}{32},\frac{23}{32},\frac9{32}\right).
\]
Here $H=9/256$, $R^2=1/256$, and the weighted guard sum is $-10/256$.
The endpoint condition is not being claimed to hold by itself at this
rational tuple. Only the full finalized equation vanishes.

\section{Consequences for unbounded computation}\label{sdf:sec:undecidable}

\subsection{An undecidable smooth quartic class}

The MRDP theorem identifies computably enumerable relations with
Diophantine relations. The existence of a fully mechanized proof, including
both directions and machine translations, is established in the work of
Larchey-Wendling and Forster \cite{coq}. Here MRDP is a background theorem,
not a result reproved by the finite counter compiler.

\begin{theorem}[A geometrically restricted complete problem]\label{sdf:thm:hardness}
Integer solvability is computably enumerable complete, under computable
many-one reductions, even on the explicit class of quartic polynomials
produced by \eqref{sdf:eq:F}. Every instance in this class comes with
\begin{enumerate}[label=\textnormal{(\alph*)}]
\item a smooth affine model over $\Z$ with geometrically integral fibres;
\item a bounded-degree integral Jacobian certificate;
\item a constructible nonnegative dyadic rational point;
\item integral points over every $\Z_p$ and solutions modulo every integer;
\item a Zariski-dense set of rational points and exactly two real path
components.
\end{enumerate}
The analogous statement for natural solvability holds on the natural
quadratic-lifting version of the class.
\end{theorem}
\begin{proof}
Given an integer polynomial equation, apply the integer version of
Lemma~\ref{sdf:lem:lifting}, then the finalizer. Integer solvability is
preserved at both stages. Equivalently, to start from a natural MRDP
representation, introduce integer variables for each original natural
witness and impose that it is a sum of four integer squares. These
additional equations are quadratic, and circuit lifting deals with the
remaining polynomial equations. This second route does not preserve
witness uniqueness at the four-square stage, and no such assertion is
needed for undecidability.

The known computably enumerable completeness of Diophantine solvability
therefore transfers to the displayed subclass. Membership remains
computably enumerable by enumerating integer tuples and evaluating the
polynomial. All geometric and arithmetic promises follow from
Theorem~\ref{sdf:thm:main}; their certificates can be generated with the
reduction and do not require discovering an integer solution.

For natural solvability, use the natural circuit lifting in
Lemma~\ref{sdf:lem:lifting}. The finalizer then preserves the full natural
zero set bijectively.
\end{proof}

This theorem is compatible with the older smooth-variety undecidability
result cited in the introduction. Its content is a specific simultaneous
normal form, including an everywhere smooth \emph{integer} model and
explicit positive evidence for all the easier existence notions.

\wtag\ [The computability core of Theorem~\ref{sdf:thm:hardness}, MRDP
followed by unique circuit lifting and a sum of squares of degree four, is
the quartic normalization of \emph{canonical-diophantine-certificates},
Proposition \path{cdc:of:prop:quartic}; the new content is the package
(a)--(e). In the formal project beside the collection, MRDP is the Lean
theorem \texttt{Diophantine.mrdp}
(\path{Computability/HilbertTenthProblem/Lean/Diophantine/MRDP.lean}), and the
natural-to-integer four-square substitution of the second route appears as
\texttt{exists\_sqTerm\_eq} in
\path{Computability/HilbertTenthProblem/Lean/Diophantine/Paper1978/Enumeration.lean};
neither formalizes any statement of this report.]

\subsection{Fixed arity versus external duration}

Fixing a universal Diophantine source relation gives a fixed finite
quadratic circuit and therefore a fixed number of variables in the
finalized family. This is a legitimate fixed-arity unbounded statement,
because the source already comes from MRDP or another proved unbounded
representation.

By contrast, the counter frontend of Section~\ref{sdf:sec:counter} creates
$n_T$ variables and $m_T$ equations. Its horizon is an external compiler
argument. The facts that each accepting bounded run has one witness and
that finalization adds only five variables do not combine to give a
fixed-arity unique representation of arbitrary halting.

More explicitly, even the disjoint union of the bounded accepting
certificates has a copy for every $T$ after the first halting time because
of padding. Restricting to a first-halting horizon would remove that
particular repetition, but compiling a variable-length history into one
fixed polynomial while preserving all multiplicities is an additional
problem. It is not accomplished by treating a loop in the compiler as a
polynomial variable.

\subsection{Input parameters and total degree}\label{sdf:sec:parameter-degree}

Suppose a master family uses an input parameter $a$ only through a loader
$v-a=0$, while the remaining quadratic equations have fixed integer
coefficients. If $a$ is treated as a coefficient during homogenization,
the loader becomes
\[
 tv-at^2.
\]
Its square has total degree six when $a$ is counted as a coordinate,
although it is quartic in witness variables for every fixed $a$.
Therefore this straightforward master presentation has
\[
 \deg_{\text{all coordinates}}\mathcal F\le6,
 \qquad \deg_{\text{witnesses}}\mathcal F(a;\cdot)=4.
\]
It is smooth \emph{relative to the input base}: the same Jacobian identity
uses only witness derivatives and continues to hold with $a$ in the
coefficient ring.

Alternatively, homogenizing $a$ as a source variable keeps the total
polynomial quartic and gives a smooth total space. It does not by itself
prove that every fixed-$a$ fibre is smooth, because the Euler identity
then uses the derivative in the $a$ direction. These two degree/smoothness
contracts must not be conflated.

\subsection{A separation between rational and integer search}

\begin{corollary}[No computable bound for the least integer witness]\label{sdf:cor:nobound}
There is no computable function of the finite source description that
bounds the height of an integer point on every solvable compiled instance.
This remains true even when the dyadic point and Jacobian certificate are
included as additional input data.
\end{corollary}
\begin{proof}
Such a function would turn integer solvability into a finite search over a
box of its computed radius, deciding the complete problem in
Theorem~\ref{sdf:thm:hardness}. The additional data are computably produced
from the same source description and do not change the reduction.
\end{proof}

The rational side is starkly different: a point of polynomially bounded
coordinate height in $\sqrt{c+1}$ is always available. Thus small rational
solutions do not bound the size of the first integer solution. The exact
height transfer in Corollary~\ref{sdf:cor:height} also shows that this
separation is not caused by huge integer auxiliary values introduced by
the finalizer.

\section{A direct higher-degree extension}\label{sdf:sec:power-two}

Quadratic lifting is useful for degree four, but it is not needed for
smoothness itself. A power-of-two padding gives a second tradeoff: keep
only five new variables and permit a larger degree.

\begin{theorem}[Power-of-two homogeneous padding]\label{sdf:thm:padding}
Let $f_i\in\Z[x]$ have degrees at most $D$. Choose a positive power of
two $d\ge\max\{1,D\}$ and set
\[
 q_i(x,t)=t^d f_i(x/t),\qquad H=\sum_iq_i^2,
 \qquad F=H+(tu-1)^2+\guard(z_1)+\guard(z_2)+2\guard(z_3),
\]
where homogenization is interpreted polynomially. Then the natural and
integer zero-set correspondences of Theorem~\ref{sdf:thm:main} still hold.
The affine model is smooth over $\Z$ with geometrically integral fibres,
and it has dyadic points and integral points at all primes. Its degree is
at most $\max\{4,2d\}$.
\end{theorem}
\begin{proof}
Nonnegativity, the unit equation, and
$q_i(\varepsilon a,\varepsilon)=\varepsilon^d f_i(a)$ prove the same
lattice correspondences. The rank-three argument in the guard variables
proves geometric integrality, and the characteristic-two fibre is again
a polynomial graph.

For smoothness, work in the quotient ring by $F$ and all its derivatives.
The equation $D_1=4z_1-1=0$ makes two invertible. The other guard
coefficients are one and two, hence are also units in this quotient, and
all $z_j=1/4$. Euler's identity now gives
\[
 \sum_i x_iF_{x_i}+tF_t-uF_u=2dH=0.
\]
Since $d$ is a power of two, $H=0$. The equation $F=0$ therefore gives
$R^2=1/2$. The derivative $F_u=2tR=0$, multiplied by $u$, gives
$R(R+1)=0$, so $R=-1/2$. Squaring contradicts $R^2=1/2$ by forcing
$1/4=0$. Since four is a unit, the quotient is the zero ring. This is
exactly the unit-Jacobian-ideal criterion for smoothness.

For dyadic points choose a power of two $K$ with $K^{2d}>2c$ and represent
\[
 K^{2d}-2c=A_1^2+A_2^2+2A_3^2+2B^2.
\]
Take $x=0$, $t=1/K$, $R=B/K^d$, $u=K+B/K^{d-1}$, and
$z_j=1/4+A_j/(2K^d)$. The same substitution as before proves $F=0$.
This point is integral at odd primes; at two the equation
$c+2z_1^2-z_1=0$ again gives a lift with $t=u=1$.
\end{proof}

Choosing the smallest such $d$ gives $d<2D$ for $D\ge1$, except that
$d=D$ when $D$ already is a power of two. Thus there is a simple
variable/degree tradeoff between direct padding and unique quadratic
circuit lifting. Only the quadratic version is implemented in the supplied
reference compiler; this section is a proved extension, not an additional
software feature.

\section{Implementation, exact checks, and reproducibility}\label{sdf:sec:implementation}

\subsection{Files and interfaces}

The archive contains a self-contained \LaTeX\ source and compiled PDF,
three Python modules, two expanded polynomial exports, and an execution
report. The main software files are
\begin{center}
\begin{tabular}{p{0.34\textwidth}p{0.56\textwidth}}
\toprule
File & Purpose\\
\midrule
\texttt{code/smooth\_compiler.py} & Validated quadratic input, homogeneous lift,
quartic export, Jacobian certificate, dyadic point, and compatible
$2$-adic residues.\\
\texttt{code/counter\_frontend.py} & Counter-program semantics and the bounded
quadratic frontend.\\
\texttt{code/verify.py} & Exact symbolic identities, independent sparse
polynomial evaluation, finite exhaustive checks, and export generation.\\
\bottomrule
\end{tabular}
\end{center}

\wtag\ [Shipped names. In this report's directory the delivered files carry
the prefix \path{13-smooth-quartic-}: \path{code/smooth_compiler.py},
\path{code/counter_frontend.py}, \path{code/verify.py} and the
\path{Makefile} are \path{code/13-smooth-quartic-*}; the two exports
\path{examples/countdown_T3.json} and \path{examples/inconsistent.json}, the
records \path{verification/results.json} and
\path{verification/document_qa.json}, and \path{requirements.txt} are
\path{data/13-smooth-quartic-*}; \path{RESEARCH_STATUS.md} and
\path{SOURCE_AUDIT.md} are at the report root under the same prefix. The
delivered PDF, the run log (a copy of \path{results.json}) and the checksum
file are not shipped; the PDF of this report is a new build. Delivered names
below are the source's. The programs import each other by their delivered
names and write into \path{examples/} and \path{verification/} beside
\path{code/}, so they must be run in a copy with the delivered layout; the
report's README gives the recipe.]

A \texttt{QuadraticSystem} copies its variables and residuals into immutable
tuples. It rejects undeclared symbols, repeated variable names,
nonintegral coefficients, floating-point coefficients, and residual degree
above two. The finalizer checks auxiliary-name collisions. Counter-program
validation checks register ranges, state targets, instruction fields,
nonnegative initial counters, and an exact nonnegative integer horizon.

The dyadic constructor computes a four-square decomposition by a finite
pair table and converts it to weights $(1,1,2,2)$ using an equal-parity
pair. The pair-table routine has a declared default input cap of one
million to prevent an unexpected large allocation. That is a software
resource limit, not a mathematical exception. The caller may supply a
four-square representation for independent exact checking instead.

\subsection{Recorded checks}

The included run used Python 3.13.5 and SymPy 1.14.0 with seed 20261002.
All 21,128 recorded assertions passed. The principal groups are shown in
Table~\ref{sdf:tab:checks}; the exact full breakdown is in
\texttt{verification/results.json}.
\wtag\ [Shipped as \path{data/13-smooth-quartic-results.json}. A rerun at
placement (Python 3.14.4 and 3.13.5, SymPy 1.14.0, on a copy) and the
Hilbert's-tenth review~\cite{review13} both reproduced all 21,128 assertions
and the 40 group counts.]

\begin{table}[htbp]
\centering
\small
\begin{tabular}{p{0.70\textwidth}r}
\toprule
Check group & Assertions\\
\midrule
Integer tuples in three explicit six-dimensional boxes & 10,125\\
Natural subboxes of those same boxes & 648\\
One-step countdown assignments, exhaustively compared to semantics & 2,048\\
Ternary obstruction, all triples modulo sixteen & 4,096\\
Primitive triples modulo eight in the same obstruction & 448\\
Independent sparse-export evaluations on seeded tuples & 520\\
Odd-prime reductions of the constructed dyadic points & 260\\
Four-square and weighted-four-square replay for $0\le N\le200$ & 402\\
Four nested Jacobian/Euler identities on 27 source systems & 108\\
Expanded B\'ezout identities and coefficient-degree bounds & 54\\
Malformed-input rejection checks & 21\\
\bottomrule
\end{tabular}
\caption{Selected groups of exact checks; not a replacement for the proofs.}
\label{sdf:tab:checks}
\end{table}

The source systems used for symbolic checks include zero-variable cases,
a zero system, constant inconsistent systems, simple quadratic equations,
20 seeded two-variable systems, and the countdown certificate. The
integer-box test ranges over $x,t,u\in\{-2,-1,0,1,2\}$ and each of the
three guard variables in $\{-1,0,1\}$. There are $5^3 3^3=3375$ tuples
per source. These ranges are recorded so that a finite test is not mistaken
for an unbounded exhaustive verification.

For the counter frontend, 140 initial-input/horizon combinations are
checked across a one-counter and a two-counter program. The transitions,
endpoint, variable count, residual count, and nonnegativity checks are
separate assertions. The one-step exhaustive check independently compares
the polynomial rows with the operational transition relation, rather than
merely checking the one run returned by the compiler.

The sparse JSON evaluator multiplies integer coefficients by powers of
supplied coordinates; it does not call SymPy substitution into the original
residual list. On seeded tuples its result is compared with direct
homogeneous-residual evaluation. For the Jacobian identity, both the nested
formula \eqref{sdf:eq:unit} and the expanded coefficient formula
\eqref{sdf:eq:bezout} are checked symbolically from the actual emitted $F$.

\subsection{What the checks establish and what they do not}

The symbolic identities certify the algebraic calculations for the exported
examples and sampled systems. The paper supplies the universal derivations.
The finite checks can reveal implementation errors in lifting, signs,
weights, inactive slacks, variable order, and exported coefficients. They
do not prove the general source semantics or replace the scheme-theoretic
arguments about smoothness and geometric integrality.

No Lean or Coq process was run for this manuscript. No existing repository
axiom audit was rerun. The ordinary proofs of rational density, the number
of real components, the higher-degree padding theorem, and the general
undecidability reduction are not advertised as consequences of a Python
unit test.

To reproduce the package from its root directory:
\begin{lstlisting}
python -m pip install -r requirements.txt
python code/verify.py
pdflatex -interaction=nonstopmode -halt-on-error article.tex
pdflatex -interaction=nonstopmode -halt-on-error article.tex
pdflatex -interaction=nonstopmode -halt-on-error article.tex
\end{lstlisting}
A minimal use of the finalizer is
\begin{lstlisting}[language=Python]
import sys
sys.path.insert(0, "code")
import sympy as sp
from smooth_compiler import QuadraticSystem, compile_smooth
from smooth_compiler import dyadic_point

x, y = sp.symbols("x y")
source = QuadraticSystem((x, y), (x - y*y, y - 2))
certificate = compile_smooth(source)
integer_point = certificate.lift((4, 2))
assert certificate.evaluate(integer_point) == 0
rational_point, metadata = dyadic_point(certificate)
assert certificate.evaluate(rational_point) == 0
assert certificate.jacobian_certificate()["unit"] == 1
\end{lstlisting}
\wtag\ [Both listings assume the delivered layout; in this report's
directory neither \path{requirements.txt} nor \path{code/verify.py} exists
under that name, and \path{verify.py} rewrites the two exports and
\path{results.json} in place. Use the copy recipe of the report's README; on
this machine the interpreter is \texttt{py}.]

\section{Formalization plan and theorem boundaries}\label{sdf:sec:boundaries}

A useful formalization should mirror the actual proof boundary rather than
importing the final statement as a large opaque axiom. The algebraic core
can be separated from the geometry.

First formalize integer guard nonnegativity and its unique zero, the two
units of $\Z$, and homogeneous lifting. These yield the natural bijection,
the signed integer bijection, and the height correspondence. Their proofs
need no algebraic geometry.

Next formalize the identities \eqref{sdf:eq:Euler}, \eqref{sdf:eq:EDJ}, and
\eqref{sdf:eq:bezout} over a general commutative coefficient ring. One may
represent the certificate as a finite tuple of polynomials and check the
single ring identity. An exported, independently checked identity is a
natural target, but a theorem connecting such a checker to polynomial
semantics is needed before it becomes a trusted bridge.

The geometric layer should separately connect a unit Jacobian ideal to
relative smoothness and prove the rank-three irreducibility argument.
The characteristic-two polynomial-graph isomorphism is elementary and
should not be subsumed under a characteristic-zero smoothness theorem.
The rational/local layer then needs four-square existence, the parity
conversion, exact rational substitution, and compatible $2$-adic lifting.

Finally, the counter frontend should prove its natural one-hot lemma,
per-edge correctness including inactive slacks, and induction on the
externally fixed horizon. Its output can then feed the generic finalizer
without repeating the geometric proof for every language.

Several boundaries are structural, not minor missing optimizations:
\begin{itemize}
\item A smooth affine model is not a smooth proper model. The points at
infinity of a homogenized projective closure are outside the main theorem.
\item Dyadic solubility of this class does not decide arbitrary equations
over $\Z[1/2]$ or $\Q$. The reduction preserves integer or natural
solvability, not rational solvability.
\item Witness preservation does not establish single-fold or finite-fold
MRDP. It preserves whatever multiplicities the source already has.
\item External horizon bounds are not compressed by adding five variables.
Fixed arity is obtained only from an already fixed-arity source relation.
\item The sharp five-auxiliary lower bound proved here is restricted to the
positive weight-four guard architecture; it is not a general lower bound
for smooth Diophantine representations.
\end{itemize}

\section{Further research questions and concrete next targets}\label{sdf:sec:future}

The following questions are proposals arising from this construction.
Unless explicitly stated otherwise, they are not assertions that no answer
exists in the literature. A dedicated priority search should accompany any
attempt to publish a claimed resolution.

\begin{question}[Four auxiliary variables with the full package]\label{sdf:q:four}
Can another integer-polynomial finalizer preserve every natural zero set
bijectively, remain quartic and smooth over $\Z$, and guarantee a dyadic
point for every source while adding at most four variables?
\end{question}
Proposition~\ref{sdf:prop:two-guard} excludes the most immediate candidate,
not all candidates. A successful construction must avoid the ternary-form
obstruction rather than merely find a better algorithm for that form.
Promising search spaces include unequal guard centres, nonquadratic guards
whose total degree stays four, or a replacement for the product-unit
constraint.

\begin{question}[Other total weights]\label{sdf:q:weights}
Within \eqref{sdf:eq:Fweights}, can two positive guards of total weight sixteen
give the complete dyadic guarantee for every quadratic source?
\end{question}
The smoothness classification allows total weight sixteen. It leaves a
finite family of two-guard choices with first weight one, or a slightly
larger family after replacing that normalization by a coprime-weight
condition. The rational point problem becomes a ternary quadratic-form
problem coupled to fourth-power scaling. An exact local obstruction or a
uniform construction would materially extend the sharp result above.

\begin{question}[A totally quartic, relatively smooth universal family]\label{sdf:q:totalquartic}
Can one obtain a single polynomial of total degree four in both input
parameters and witnesses, with every integer-input fibre smooth over $\Z$,
while preserving a universal natural relation and retaining the same
local-solubility guarantees?
\end{question}
Section~\ref{sdf:sec:parameter-degree} gives a total-degree-six master family
and explains why homogenizing the parameter changes the relative-smoothness
claim. An input loader compatible with both total degree and the relative
Euler identity is the specific missing component.

\begin{question}[A smaller Jacobian certificate]\label{sdf:q:jacobian}
Can the degree-seven expanded identity be replaced uniformly by a
certificate whose products all have degree at most six, or by a shorter
arithmetic circuit with the same universal coefficient contract?
\end{question}
This is a certificate-optimization problem, not an attempt to improve a
universal halting polynomial's operation count by ignoring its frontend.
The precise optimization object should include the equation, all derivative
costs, and the coefficient tuple in \eqref{sdf:eq:coefficients}.

\begin{question}[Controlled higher topology]\label{sdf:q:topology}
Which topological invariants of the two real components can be read from
the real zero set of the source system, and can the construction be modified
to make every real component contractible?
\end{question}
The model \eqref{sdf:eq:real-model} gives a concrete starting point. Its fibres
change from a capped sphere family to a cylinder at source real zeros.
The number of components is fixed, but a full homotopy or diffeomorphism
classification is not supplied here.

\begin{question}[Strong approximation and integral obstructions]\label{sdf:q:strong}
What can be said about strong approximation, and about more refined
obstructions to integral points, on these explicit affine varieties?
\end{question}
All local integral points exist and rational points are dense in the real
locus, but these statements do not imply simultaneous approximation at
many places with integral control away from them. A meaningful answer must
respect the undecidability theorem: no proposed effectively complete,
uniform terminating obstruction procedure can decide all compiled inputs.

\begin{question}[Density of restricted-denominator points]\label{sdf:q:dyadicdensity}
For which source systems are the dyadic points dense in a real component?
Can a bounded set of additional denominator primes guarantee density?
\end{question}
The proof of Theorem~\ref{sdf:thm:density} uses rational parametrizations with
unrestricted denominators. Replacing them by dyadic parameters does not
preserve all denominator conditions automatically. This is a distinct
arithmetic approximation problem, not an immediate corollary of the
existence construction.

\begin{question}[Smooth models over other arithmetic rings]\label{sdf:q:rings}
Which discretely ordered domains or rings of $S$-integers admit analogues
with exact control of unit-induced witness multiplicity?
\end{question}
The integer proof uses both the positivity gap of $2z^2-z$ on the lattice
and the fact that the unit group is $\{\pm1\}$. Neither property carries
over unchanged to a general ring of integers. A useful extension should
state exactly what replaces the lattice order and how the unit group acts
on fibres.

\begin{question}[A verified finalizer integrated with repository frontends]\label{sdf:q:verified}
Can the five-variable transformation be formalized once and instantiated
on the repository's existing causal, queue, heap, and interaction-net
interfaces without changing their semantic contracts?
\end{question}
A practical first milestone is the homogeneous lift plus the explicit
Jacobian certificate. A second is a proved checker for the sparse JSON
format. A third is one full frontend-to-smooth-quartic theorem with an
explicit distinction between natural and integer witnesses.
\wtag\ [Open at the write. The interfaces named are Parts of the
collection's report \emph{canonical-diophantine-certificates}, which are
research text, not formal developments; the Lean project
\path{Computability/HilbertTenthProblem} has no smoothness or finalizer
module, and the review's portable helper~\cite{review13} is an independent
replay, not a proved checker.]

\begin{question}[A detailed publication and priority audit]\label{sdf:q:priority}
Which parts of the combined normal form are already implicit in known
smoothness reductions or in work on integral local-to-global phenomena?
\end{question}
The introduction identifies a definite antecedent and avoids presenting
the general smooth-undecidability phenomenon as new. A systematic comparison
should separately track good reduction at every prime, fixed auxiliary
count, dyadic sections after specialization, the exact natural bijection,
and the bounded-degree Jacobian certificate. These properties should not
be collapsed into the vague phrase ``smooth Diophantine representation.''
\wtag\ [Open at the write. The Hilbert's-tenth review~\cite{review13} checked
Poonen's Lemma~4.1 and Corollary~4.3, the mechanized MRDP record and the
Stacks Project criterion, and states that it was not a complete historical
priority search.]

\section{Conclusion}\label{sdf:sec:conclusion}

The construction separates three layers that are easily confused. The source
quadratic equations contain the computational semantics. The integer gap of
the guards and the unit equation preserve their witnesses. The homogeneous
scaling and weighted guards alter the surrounding arithmetic geometry so
that every geometric fibre is smooth and integral, all integral local
points exist, and rational points are plentiful.

The resulting polynomial is
\[
 \boxed{\sum_i q_i(x,t)^2+(tu-1)^2+
       (2z_1^2-z_1)+(2z_2^2-z_2)+2(2z_3^2-z_3).}
\]
It has exactly five additional variables. Its geometric regularity is
certified by an integral identity of bounded degree, not by a generic
choice, numerical evidence, or an unbounded singularity search. Its
natural witnesses are unchanged; its integer witnesses form exactly two
sign-related copies. The failed two-guard alternative shows that the
arithmetic point construction imposes a genuine additional constraint
beyond smoothness.

This does not settle a classical undecidability conjecture or erase the
finite-fold boundary. It provides a reusable explicit normal form: even a
smooth affine quartic with everywhere good affine reduction, small dyadic
points, and only two real components can carry the full difficulty of
integer computation.

\appendix
\section{The complete countdown residual list}\label{sdf:app:countdown}

For the input two and horizon $T=3$, use variables
\[
 q_{s,0},q_{s,1},c_s\quad(0\le s\le3),\qquad
 a_{s,D},a_{s,Z},a_{s,H},d_s\quad(0\le s<3).
\]
There are $3\cdot4+4\cdot3=24$ of them. The initial rows are
\[
 q_{0,0}-1,\qquad q_{0,1},\qquad c_0-2.
\]
At each $s=0,1,2$, include the following eight residuals:
\begin{align*}
 &a_{s,D}+a_{s,Z}+a_{s,H}-1,\\
 &q_{s,0}-a_{s,D}-a_{s,Z},\qquad q_{s,1}-a_{s,H},\\
 &q_{s+1,0}-a_{s,D},\qquad q_{s+1,1}-a_{s,Z}-a_{s,H},\\
 &c_{s+1}-c_s+a_{s,D},\\
 &a_{s,D}c_s-a_{s,D}-d_s,\qquad a_{s,Z}c_s.
\end{align*}
The final row is $q_{3,1}-1$. Thus the total number is $3+3\cdot8+1=28$.
The source witness is
\begin{align*}
 &(q_{0,0},q_{0,1})=(q_{1,0},q_{1,1})=(q_{2,0},q_{2,1})=(1,0),\\
 &(q_{3,0},q_{3,1})=(0,1),\qquad (c_0,c_1,c_2,c_3)=(2,1,0,0),\\
 &(a_{0,D},a_{0,Z},a_{0,H})=(a_{1,D},a_{1,Z},a_{1,H})=(1,0,0),\\
 &(a_{2,D},a_{2,Z},a_{2,H})=(0,1,0),\qquad(d_0,d_1,d_2)=(1,0,0).
\end{align*}
Add $(t,u,z_1,z_2,z_3)=(1,1,0,0,0)$ to obtain the unique natural zero of
the finalized equation.

For an entirely explicit compact polynomial, apply the following rule to
each row above: multiply its quadratic terms by $1$, its linear terms by
$t$, and its constants by $t^2$. Square the 28 resulting homogeneous
polynomials, add them, and then add $(tu-1)^2+\guard(z_1)+\guard(z_2)+
2\guard(z_3)$. This is a complete expression, not an appeal to an
unspecified universal representation. Its expanded 97-term coefficient
list is in \texttt{examples/countdown\_T3.json}.

\section{The export and verification boundary}\label{sdf:app:export}

Each JSON export records an ordered variable list and a list of terms.
A term has an integer coefficient and an exponent vector of the same
length as the variable list. Its mathematical meaning is
\[
 c\prod_{j=1}^Nv_j^{e_j}.
\]
Zero coefficients are omitted; the degree is the maximum sum of exponents
among retained terms. The format also records the source residuals for
human inspection, the guard weights, a rational point as exact strings,
the four-square conversion data, and compatible $2$-adic residues.

The equation is reconstructed for independent evaluation from the
coefficient/exponent data. The source residual strings are not interpreted
as trusted assertions about machine semantics. A future formal checker
should validate the term format, coefficient ring, variable order, and
identity semantics before using such an export as a theorem input.

The file \texttt{examples/inconsistent.json} uses the constant source
residual $1$ with one unused source variable. It therefore has one more
free source coordinate than Example~\ref{sdf:ex:one}; it is the product of
that example with an affine line. The printed five-variable equation and
its rational point are checked separately by the zero-variable test.
This distinction is intentional and is not a mismatch in the auxiliary
variable count.

\wtag\ [The two exports are shipped as
\path{data/13-smooth-quartic-countdown_T3.json} and
\path{data/13-smooth-quartic-inconsistent.json}.]

\section{Provenance and scope of the source review}\label{sdf:app:sourcereview}

The source comparison consisted of the repository root overview, the
Hilbert-tenth-problem overview, the canonical-certificate collection
README, and a pinned canonical-certificate research-status file. The
last file was read at commit
\begin{center}\small\ttfamily a845ab5d0f4591294ef3c2d92feaeb5825a8d543\end{center}
and had Git blob identifier
\begin{center}\small\ttfamily e768b2284c6723d21d77edfb24f786ccdb83585a.\end{center}
The source review did not establish the absence of an equivalent formula
elsewhere in the repository or literature.

The external mathematical references used directly are Poonen's smoothness
reduction, the primary account of mechanized MRDP by Larchey-Wendling and
Forster, and the Stacks Project's smoothness criterion. The general
four-square theorem is used as classical background; its application to
weighted $(1,1,2,2)$ representation is proved in the text. The explicit
five-variable formula, the universal B\'ezout identity, the dyadic
construction, the weighted two-guard obstruction, and the detailed
counter frontend were developed in this manuscript and checked as
reported. Their status is that of proved research claims awaiting
independent mathematical review, not established publication priority.

\section{Provenance of this report}\label{sdf:app:provenance}

\wtag\ The report was built from \emph{one} manuscript, batch-79
manuscript~13 of ProveIt's incoming reports (cluster J1 of that batch). It
arrived as the archive \path{Smooth_Quartic_Diophantine_Certificates.zip} in
commit \texttt{ef2fc7990} (2 October 2026), with inner directory
\path{smooth_diophantine/}: the main file \path{article.tex} (1,733 lines,
printed here), a 28-page PDF, a delivery README, research-status and
source-audit notes, three Python modules, a Makefile, two exports, two
verification records, a run log, a requirements file and a checksum file
(16 members; the checksum file verified 15 of 15 at placement). Commit
\texttt{224ca41df} placed it as this new report and staged its programs,
data and notes byte-identical to the delivery.

\paragraph{Pin.} The manuscript names the repository commit
\path{a845ab5d0f4591294ef3c2d92feaeb5825a8d543} and the blob
\path{e768b2284c6723d21d77edfb24f786ccdb83585a} of
\path{canonical-diophantine-certificates/01-causal-traces-RESEARCH_STATUS.md}
(Section~\ref{sdf:sec:repo} and Appendix~\ref{sdf:app:sourcereview}); its
source audit also names the root README blob \texttt{bc4c514b61e3}. Both
blobs are unchanged at the write.

\paragraph{Contribution.} Everything in the report except the text marked
\wtag: the finalizer and Theorem~\ref{sdf:thm:main} with its proofs, the
size bounds, the weighted-guard classification and the restricted
optimality, the counter frontend and its worked example, the undecidability
corollaries, the power-of-two padding, the implementation and its recorded
checks, ten research questions, and two appendices on the countdown
certificate and the export boundary.

\paragraph{Where the write had to choose.}
\begin{itemize}
\item \emph{Placement.} The manuscript post-processes any quadratic system,
CDC's or not, and its subject is the arithmetic geometry of the final
equation rather than a substrate, so it was placed as a new report rather
than as a further Part of \emph{canonical-diophantine-certificates} after
that report's Part~XI (the alternative recorded at placement).
\item \emph{Re-proved results.} Lemma~\ref{sdf:lem:lifting},
Theorem~\ref{sdf:thm:counter} and the computability core of
Theorem~\ref{sdf:thm:hardness} re-present results of that report. They are
kept in full, because later proofs use them, with \wtag\ pointers naming
the host labels and no novelty claimed.
\item \emph{Numbering.} The editorial preface is unnumbered, so the
manuscript's section, statement, equation and table numbers are unchanged;
this appendix is Appendix~\ref{sdf:app:provenance}, after the manuscript's
Appendices~\ref{sdf:app:countdown}--\ref{sdf:app:sourcereview}.
\item \emph{Text kept as delivered.} The title page, the author line
(``Prepared for Vladimir Reshetnikov / Developed with ChatGPT'') and the PDF
metadata are the source's; one \wtag\ sentence was added to the status box.
The display ``$w_j>0$'' of Corollary~\ref{sdf:cor:five} is kept, with a
\wtag\ note giving $w_j\in\Z_{>0}$. No symbol was renamed and no
normalization changed; the preface's notation table lists the reused
letters.
\item \emph{Added text.} The preface; \wtag\ notes in
Sections~\ref{sdf:sec:repo}, \ref{sdf:sec:novelty}, \ref{sdf:sec:interface}
(after Lemma~\ref{sdf:lem:lifting}), \ref{sdf:sec:local} (after
Lemma~\ref{sdf:lem:1122}), \ref{sdf:sec:guards} (after
Corollary~\ref{sdf:cor:five}), \ref{sdf:sec:counter},
\ref{sdf:sec:undecidable}, \ref{sdf:sec:implementation} (shipped file names,
reruns) and \ref{sdf:sec:future} (Questions~\ref{sdf:q:verified}
and~\ref{sdf:q:priority}) and in Appendix~\ref{sdf:app:export}; this
appendix; and the bibliography entries
\cite{cdcreport,pqcreport,review13,ramanujan,minsky}.
\end{itemize}

\paragraph{Files.} The report directory holds this text
(\path{article.tex}, \path{article.pdf}), a report README, the delivered
\path{RESEARCH_STATUS.md} and \path{SOURCE_AUDIT.md} as
\path{13-smooth-quartic-RESEARCH_STATUS.md} and
\path{13-smooth-quartic-SOURCE_AUDIT.md}, the programs and Makefile in
\path{code/} and the exports, records and requirements file in
\path{data/}, all under the prefix \path{13-smooth-quartic-}. Not shipped:
the delivered PDF and README, the run log (the results record plus one
newline) and the checksum file. They survive in the archive of commit
\texttt{ef2fc7990}.

\begin{thebibliography}{9}

\bibitem{poonen}
Bjorn Poonen, with an appendix by Matthias Aschenbrenner,
\emph{Automorphisms mapping a point into a subvariety},
Journal of Algebraic Geometry \textbf{20} (2011), 785--794.
Author manuscript dated 14 July 2009, especially Lemma 4.1,
Corollary 4.3, and Section 5.
\url{https://math.mit.edu/~poonen/papers/automorphism.pdf}.

\bibitem{coq}
Dominique Larchey-Wendling and Yannick Forster,
\emph{Hilbert's Tenth Problem in Coq (Extended Version)},
Logical Methods in Computer Science \textbf{18}(1), article 35 (2022).
\url{https://arxiv.org/abs/2003.04604}.

\bibitem{stacks}
The Stacks Project Authors,
\emph{The Stacks Project}, Tag 01V4, Smooth morphisms.
\url{https://stacks.math.columbia.edu/tag/01V4}.
Accessed 2 October 2026.

\bibitem{repo}
Vladimir Reshetnikov and contributors,
\emph{ProveIt}, repository and Hilbert-tenth-problem overview.
\url{https://github.com/VladimirReshetnikov/ProveIt}.
Repository paths \texttt{README.md} and
\texttt{Computability/HilbertTenthProblem/README.md}; accessed
2 October 2026.

\bibitem{canonical}
ProveIt research collection,
\emph{Canonical Diophantine certificates}, collection README.
Repository path
\path{SetTheory/Cardinals/docs/reports/hilbert-tenth-problem/canonical-diophantine-certificates/README.md}.
Accessed 2 October 2026. Research-report claims are not treated as
proof-assistant verification in this manuscript.

\bibitem{status}
ProveIt research collection,
\emph{Research and verification status}, canonical causal certificates.
Repository path
\path{SetTheory/Cardinals/docs/reports/hilbert-tenth-problem/canonical-diophantine-certificates/01-causal-traces-RESEARCH_STATUS.md}.
Read at commit \texttt{a845ab5d0f4591294ef3c2d92feaeb5825a8d543};
blob \texttt{e768b2284c6723d21d77edfb24f786ccdb83585a}.

\bibitem{cdcreport}
\wtag\ ProveIt research-report collection,
\emph{Canonical Diophantine certificates}, report directory
\path{SetTheory/Cardinals/docs/reports/hilbert-tenth-problem/canonical-diophantine-certificates/};
cited by label (\path{cdc:}\dots), which does not resolve in this file.
AI-assisted, unrefereed, not formalized.

\bibitem{pqcreport}
\wtag\ ProveIt research-report collection,
\emph{Probabilistic, quantum and continuous computation}, report directory
\path{SetTheory/Cardinals/docs/reports/hilbert-tenth-problem/probabilistic-quantum-and-continuous-computation/};
cited by label (\path{pqc:}\dots). AI-assisted, unrefereed, not
formalized.

\bibitem{review13}
\wtag\ ProveIt, Hilbert's-tenth research tree,
\emph{Review: Smooth Quartic Diophantine Certificates},
\path{Computability/HilbertTenthProblem/Papers/research-wip/native-stream-queue/review_smooth_quartic_aebfa386e.md},
with replay helper and receipt (\texttt{.py}, \texttt{.json}), commit
\texttt{fd4a2e8d3}, 2 October 2026.

\bibitem{ramanujan}
\wtag\ S. Ramanujan,
\emph{On the expression of a number in the form $ax^2+by^2+cz^2+du^2$},
Proceedings of the Cambridge Philosophical Society \textbf{19} (1917),
11--21.

\bibitem{minsky}
\wtag\ M. L. Minsky,
\emph{Computation: Finite and Infinite Machines},
Prentice-Hall, Englewood Cliffs, NJ, 1967.

\end{thebibliography}
\end{document}
